\documentclass[10pt,reqno]{amsart}
\usepackage[T1]{fontenc}
\usepackage{lmodern}
\usepackage{microtype}
\usepackage{amsmath,amssymb,mathtools,mathrsfs}
\usepackage{enumitem}
\usepackage[colorlinks=true,linkcolor=blue,citecolor=blue,urlcolor=blue]{hyperref}

\newtheorem{theorem}{Theorem}[section]
\newtheorem{lemma}[theorem]{Lemma}
\newtheorem{proposition}[theorem]{Proposition}
\newtheorem{corollary}[theorem]{Corollary}
\newtheorem{hypothesis}[theorem]{Hypothesis}
\newtheorem{remark}[theorem]{Remark}
\theoremstyle{definition}
\newtheorem{definition}[theorem]{Definition}

\newcommand{\Rcal}{\mathscr R}
\newcommand{\Av}{\operatorname {Av}}
\newcommand{\rad}{\operatorname {rad}}
\newcommand{\li}{\operatorname {li}}
\newcommand{\lcm}{\operatorname {lcm}}
\newcommand{\e}{\mathrm e}
\newcommand{\1}{\mathbf 1}
\newcommand{\eps}{\varepsilon}
\newcommand{\PF}{H_{\mathrm{PF}}}
\newcommand{\PFp}{H_{\mathrm{PF}}'}

% TODO(submission): replace anonymous author metadata after authorship is settled.
\title[Intrinsic congruences and a provisional $3/4$ bound]{The complete intrinsic congruence system for the Erd\H{o}s--Straus conjecture:\\ exact structure and a provisional improvement of the exceptional-set exponent}
\author{Anonymous}
\date{Draft, August 2026}

\begin{document}
\begin{abstract}
Subject to external verification, we obtain the unconditional bound
\[
 E_{\rm all}(N)\ll N\exp\{-c(\log N)^{3/4}\}
\]
for Erd\H{o}s--Straus exceptional denominators.  This theorem is explicitly \textsc{claimed}/\allowbreak\textsc{provisional}; its load-bearing pruned cubic prime-slice input, Theorem~\ref{thm:pruned} (the analogue of Theorem~34.8 in the source notes), is likewise \textsc{claimed}/\allowbreak\textsc{provisional}.  No assertion in this draft should be cited as an established improvement before independent expert review.

For each $M\equiv3\pmod4$, put $A=(M+1)/4$.  The congruence identities studied here give the intrinsic forbidden set
\[
 \Rcal(M)=\{-4D\pmod M:D\mid A^2\}.
\]
We prove this description with honest deduplication and show that its total class mass through $X$ is of exact order $(\log X)^3$.  The prime-modulus subsystem has mass $\Theta((\log X)^2)$ and avoider density $\exp\{-\Theta((\log X)^2)\}$.  Every square avoids the full intrinsic system.

Separately, a weighted rough-modulus estimate and the Lov\'asz local lemma refute the literal all-range partition-free hypothesis $\PF$ by proving
\[
 \delta_X\geq\exp\!\left\{-O\!\left(\frac{(\log X)^3\log\log\log X}{\log\log X}\right)\right\}
       =\exp\{-o((\log X)^3)\}.
\]
That refutation does not apply in the critical window $\log N\asymp(\log X)^4$.  The provisional $3/4$ argument instead uses a fixed $c$-free atom family, residue-resolved factorial moments, a uniform conditioned void supplied by Theorem~\ref{thm:pruned}, and an even-Bonferroni majorant with the honest coefficient ledger $\exp\{O((\log X)^4)\}$.  The literal $\exp\{O(J)\}$ ledger clause of the earlier critical-window hypothesis is not proved and is deliberately bypassed.  At the separate residue-dispersion frontier, the regular progression tail is proved negligible; the short-cofactor endpoint remains \textsc{open}.
\end{abstract}

\maketitle
\tableofcontents

\section{Introduction and statement map}
The Erd\H{o}s--Straus conjecture asks whether, for every integer $n\geq2$, there are positive integers $x,y,z$ such that
\begin{equation}\label{eq:ES}
 \frac4n=\frac1x+\frac1y+\frac1z.
\end{equation}
It remains open.  The solution parametrizations and average counts of Elsholtz--Tao \cite{ElsholtzTao}, the divisor criterion presented by Bloom--Elsholtz \cite{BloomElsholtz}, and the exceptional-set argument reconstructed by Pomerance--Weingartner \cite{PomeranceWeingartner} form the nearest literature context.  Yamamoto's reciprocity analysis \cite{Yamamoto}, Salez's computation \cite{Salez}, and the later counting work of Elsholtz--Planitzer \cite{ElsholtzPlanitzer} provide additional context; none is used to promote a conditional or provisional statement below.

The headline assertion of this draft is
\[
 \boxed{\ E_{\rm all}(N)\ll N\exp\{-c(\log N)^{3/4}\}\ }
 \qquad\text{\textsc{claimed/provisional}.}
\]
It is unconditional in the sense that no unproved hypothesis is assumed, but it is subject to external verification.  The load-bearing prime-slice input, Theorem~\ref{thm:pruned}, has the same provisional status.

The paper has five layers.
\begin{enumerate}[leftmargin=*,label=(\roman*)]
\item Section~\ref{sec:criterion} derives a complete two-case divisor criterion for prime denominators and the identity families that lead to the intrinsic congruence system.
\item Sections~\ref{sec:intrinsic}--\ref{sec:prime} determine that system exactly.  The distinct-modulus class mass is cubic, while the full prime slice and its avoidance exponent are quadratic.  Squares survive every modulus, and a multi-shift certificate gives a subcubic-cost lower bound for a growing part of the composite system.
\item Sections~\ref{sec:hpf}--\ref{sec:refutation} state the literal cubic-mass majorant hypothesis $\PF$, prove its formerly desired conditional consequence, and then refute the hypothesis with all quantifiers visible.  The analytic inputs are upper mean-value estimates of Nair--Tenenbaum \cite{NairTenenbaum} and Henriot \cite{Henriot}; Shiu's theorem \cite{Shiu} is used in the prime-harvesting arguments.
\item Section~\ref{sec:window} separates that all-range refutation from the critical window and records the cylinder, tensor, finite-transfer, and factorial-moment walls.
\item Section~\ref{sec:records} states the provisional record inputs.  Section~\ref{sec:cfree} then gives the complete fixed-$c$-free assembly leading, provisionally, to the $3/4$ exponent.  Its final subsection, Section~\ref{sec:dispersion-frontier}, records the proved regular-tail result and the open short-cofactor endpoint from the residue-dispersion frontier.  Scott--Sokal's partition-function correspondence \cite{ScottSokal} does not itself provide the required pointwise minorant.
\end{enumerate}

Labels are part of the assertions.  ``Proved'' means proved in this draft.  ``Open'' means no proof is claimed.  The results in Sections~\ref{sec:records} and~\ref{sec:cfree} are reproduced from the campaign record and remain \textsc{claimed/provisional}: writing out their internal arguments does not replace independent priority and referee checks.  In particular, no result of this paper proves the Erd\H{o}s--Straus conjecture.

The historical comparison is also provisional.  Vaughan's 1970 bound \cite{Vaughan} is the record identified by the search.  Pomerance--Weingartner's 2025 preprint \cite[\S4]{PomeranceWeingartner} cites Vaughan as the state of the art, and a fresh search found no post-1970 improvement of the exceptional-set exponent.  Access to Vaughan's primary paper was blocked, although its DOI was confirmed.  Thus, if the $3/4$ theorem survives external review, it would be the first exponent improvement found since 1970.  This is a search result, not a literature guarantee.

\section{Notation and elementary tools}\label{sec:notation}
All logarithms are natural.  We write $\tau_k(n)$ for the ordered $k$-fold divisor function, $\tau=\tau_2$, $\omega(n)$ for the number of distinct prime factors, $P^-(n)$ for the least prime factor, and $\varphi$ for Euler's function.  Constants implicit in $O,\ll,\asymp$ are absolute unless a subscript says otherwise.

For $X\geq3$, put $L=\log X$.  Once the intrinsic sets have been defined, their full avoider set is
\begin{equation}\label{eq:Av-def}
 \Av_X(N)=\{n\leq N:n\bmod M\notin\Rcal(M)
       \text{ for every }M\leq X,\ M\equiv3\pmod4\}.
\end{equation}
For fixed $X$ this predicate is periodic, so its natural density $\delta_X$ exists.  The superscript ``prime'' restricts the moduli to primes $\ell\equiv3\pmod4$.

We use three standard inputs: Bombieri--Vinogradov in a fixed power-saving form, Brun--Titchmarsh, and the quantitative Lov\'asz local lemma.  The latter says that events $A$ with a dependency graph have positive joint avoidance probability at least $\prod_A(1-x_A)$ if
\[
 \Pr(A)\leq x_A\prod_{B\sim A}(1-x_B),\qquad 0\leq x_A<1.
\]
The product-space form and its relation to the independent-set polynomial are discussed in \cite{ScottSokal}.

\section{Criterion and identity families}\label{sec:criterion}

\subsection{A two-term lemma and the complete prime criterion}
\begin{lemma}[two-term criterion]\label{lem:two-term}
Let $a,b$ be coprime positive integers.  The equation $a/b=1/y+1/z$ has a solution in positive integers if and only if $b^2$ has a divisor $d\equiv-b\pmod a$.  Its solutions are
\[
 y=\frac{b+d}{a},\qquad z=\frac{b+b^2/d}{a}.
\]
\end{lemma}
\begin{proof}
If $d\mid b^2$ and $d\equiv-b\pmod a$, then $d'=b^2/d$ also satisfies $d'\equiv-b\pmod a$, since $b$ is invertible modulo $a$.  The displayed $y,z$ are integral and
\[
 (d+b)(d'+b)=b(d+d'+2b),
\]
which gives $1/y+1/z=a/b$.  Conversely, from $ayz=b(y+z)$ put $d=ay-b$ and $d'=az-b$.  Positivity follows from $1/y<a/b$ and $1/z<a/b$, and
$dd'=a^2yz-ab(y+z)+b^2=b^2$.  Thus $d\mid b^2$ and $d\equiv-b\pmod a$.
\end{proof}

If $p$ is an odd prime in a solution of \eqref{eq:ES}, at least one denominator is divisible by $p$, because $4xyz=p(xy+yz+zx)$, and not all three are, since after division their reciprocal sum would be at most $3$.  Thus exactly one or exactly two denominators are divisible by $p$.

\begin{theorem}[complete divisor criterion]\label{thm:criterion}
Let $p$ be an odd prime.  Equation \eqref{eq:ES} is soluble if and only if at least one of the following holds.
\begin{description}[leftmargin=2.4em,style=nextline]
\item[(B)] There is $q>0$ with $q\equiv-p\pmod4$ and a divisor $d\mid x_0^2$, where $x_0=(p+q)/4$, such that $q\mid d+x_0$.
\item[(A)] There is $m>0$ with $m\equiv3p\pmod4$ and a divisor $d\mid z_0^2$, where $z_0=(pm+1)/4$, such that $m\mid d+z_0$.
\end{description}
The corresponding solutions are respectively
\[
 \left(x_0,\frac{p(x_0+d)}q,\frac{p(x_0+x_0^2/d)}q\right),
 \qquad
 \left(\frac{z_0+d}m,\frac{z_0+z_0^2/d}m,pz_0\right).
\]
\end{theorem}
\begin{proof}
If two denominators are divisible by $p$, write them $py',pz'$.  Then
\[
 \frac4p-\frac1{x_0}=\frac1p\left(\frac1{y'}+\frac1{z'}\right),
 \qquad q:=4x_0-p>0,
\]
so $q/x_0=1/y'+1/z'$.  Also $(q,x_0)=(p,x_0)=1$, since $p\mid x_0$ would put $p$ in all three denominators.  Lemma~\ref{lem:two-term} gives (B), both formulas, and the converse.

If exactly one denominator is divisible by $p$, call it $pz_0$.  The equation becomes
\[
 xy(4z_0-1)=pz_0(x+y).
\]
Since $p\nmid xy$, write $4z_0-1=pm$.  Then $m/z_0=1/x+1/y$ and $(m,z_0)=1$.  Lemma~\ref{lem:two-term} gives (A), while $m\equiv3p\pmod4$ is exactly the integrality condition for $z_0$.  Reversing the calculation proves the converse.
\end{proof}
This is an independent derivation of the criterion whose shape agrees with Theorem~1 of Bloom--Elsholtz \cite{BloomElsholtz}.

\subsection{Forced identities}
For $D=\prod r^{e_r}$ define
\begin{equation}\label{eq:R0}
 R_0(D)=\prod_r r^{\lceil e_r/2\rceil}.
\end{equation}
Thus $R_0(D)$ is the least positive integer $R$ for which $R\mid x$ forces $D\mid x^2$.

Fix $q,d$ in case (B) and an $R$ divisible by $R_0(d)$.  The simultaneous conditions
\[
 x_0\equiv0\pmod R,\qquad x_0\equiv-d\pmod q,
 \qquad p=4x_0-q
\]
form, when consistent, one progression modulo $4\lcm(R,q)$, and every member is represented by Theorem~\ref{thm:criterion}.  The case-(A) construction is identical with $m$ in place of $q$.  Two classical-looking forms are also immediate.

\begin{proposition}[elementary identity families]\label{prop:families}
For positive integers $g,u,v$,
\[
 p=4guv-u-v\quad\Longrightarrow\quad
 \frac4p=\frac1{guv}+\frac1{pgu}+\frac1{pgv}.
\]
For $j\geq1$, the choice $q=(x_0+1)/j$ and $d=1$ in case (B) gives
\[
 jp=(4j-1)x_0-1,
\]
and hence covers exactly the class $p\equiv-4\pmod{4j-1}$.
\end{proposition}
\begin{proof}
The first right side has common denominator $pguv$ and numerator $p+u+v=4guv$.  For the second, substitute $q=(x_0+1)/j$ into $p=4x_0-q$.  Since $4j\equiv1\pmod{4j-1}$, integrality is equivalent to $p\equiv-4\pmod{4j-1}$; then $q\mid x_0+1=d+x_0$ and Theorem~\ref{thm:criterion} applies.
\end{proof}

The identity from which the complete intrinsic system arises is the following multiplier form.
\begin{lemma}[intrinsic multiplier identity]\label{lem:multiplier}
Let $M\equiv3\pmod4$, $A=(M+1)/4$, and $A=uvw$ a factorization into positive integers.  If the positive integer $n$ satisfies $nv\equiv-u\pmod M$ and $s=(nv+u)/M$, then
\begin{equation}\label{eq:mult-id}
 \frac4n=\frac1{suw}+\frac1{nsvw}+\frac1{nuvw}.
\end{equation}
In particular $n\equiv-uv^{-1}\pmod M$ is a forced-solubility class.
\end{lemma}
\begin{proof}
Since $(A,M)=1$, the factor $v$ is invertible modulo $M$, and $s$ is a positive integer.  On the common denominator $nsuvw$, the numerator on the right is
$nv+u+s=sM+s=s(M+1)=4suvw$.
\end{proof}

\section{The complete intrinsic system}\label{sec:intrinsic}
For $M\equiv3\pmod4$, define
\[
 \Rcal(M)=\{-uv^{-1}\pmod M:uvw=(M+1)/4\},
 \qquad F(M)=|\Rcal(M)|.
\]
``Complete'' always means complete for the multiplier identity of Lemma~\ref{lem:multiplier}; it does not assert that all possible representations of \eqref{eq:ES} arise from one fixed $M$.

\begin{lemma}[exact structure and deduplication]\label{lem:exact-R}
Put $A=(M+1)/4$.  Then
\begin{equation}\label{eq:exact-R}
 \Rcal(M)=\{-4D\pmod M:D\mid A^2\},
 \qquad \frac{\tau(A^2)-1}{2}\leq F(M)\leq\tau(A^2).
\end{equation}
The lower-bound classes may be taken from the divisors $D<A$, and they are pairwise distinct.
\end{lemma}
\begin{proof}
A factorization $A=uvw$ gives $D=u^2w\mid A^2$, and $4A\equiv1\pmod M$ gives
\[
 -uv^{-1}\equiv-u^2wA^{-1}\equiv-4u^2w=-4D\pmod M.
\]
Conversely, write $A=\prod p^{e_p}$ and $D=\prod p^{d_p}$ with $0\leq d_p\leq2e_p$.  Set
\[
 v_p(u)=\lfloor d_p/2\rfloor,\quad
 v_p(w)=d_p-2\lfloor d_p/2\rfloor,\quad
 v_p(v)=e_p-v_p(u)-v_p(w).
\]
The last number is nonnegative; hence $uvw=A$ and $u^2w=D$.  This proves the set identity and upper bound.  The involution $D\mapsto A^2/D$ has unique fixed point $A$, so exactly $(\tau(A^2)-1)/2$ divisors have $D<A$.  For them $0<4D<M$, making the residues $-4D$ distinct.
\end{proof}
The set depends on $M$, not on a factorization $M=k\ell$.  Thus a sum over pairs $(k,\ell)$ must be quotiented by their product before it represents distinct intrinsic moduli.

\begin{theorem}[exact cubic supply]\label{thm:cubic}
As $Q\to\infty$,
\begin{equation}\label{eq:cubic}
 \sum_{\substack{M\leq Q\\M\equiv3\ (4)}}\frac{F(M)}M\asymp(\log Q)^3.
\end{equation}
More precisely,
\[
 \left(\frac1{48\zeta(2)}+o(1)\right)(\log Q)^3
 \leq\sum_{\substack{M\leq Q\\M\equiv3\ (4)}}\frac{F(M)}M
 \leq\left(\frac1{24\zeta(2)}+o(1)\right)(\log Q)^3.
\]
\end{theorem}
\begin{proof}
At each prime,
\[
 \sum_{e\geq0}(2e+1)z^e=\frac{1-z^2}{(1-z)^3},
\]
so
\[
 \sum_{n\geq1}\frac{\tau(n^2)}{n^s}=\frac{\zeta(s)^3}{\zeta(2s)},
 \qquad
 \tau(n^2)=\sum_{d^2\mid n}\mu(d)\tau_3(n/d^2).
\]
The elementary three-dimensional hyperbola argument gives
\[
 \sum_{m\leq x}\tau_3(m)=\tfrac12x(\log x)^2+O(x\log x).
\]
Substitution in the coefficient identity, using absolute convergence of the $d^{-2}$ sum, yields
\[
 \sum_{n\leq x}\tau(n^2)=\frac{x(\log x)^2}{2\zeta(2)}+O(x\log x).
\]
Partial summation gives
\begin{equation}\label{eq:tau-harmonic}
 \sum_{n\leq x}\frac{\tau(n^2)}n
 =\frac{(\log x)^3}{6\zeta(2)}+O((\log x)^2).
\end{equation}
Write $M=4A-1$.  Since $(4A-1)^{-1}=(4A)^{-1}+O(A^{-2})$ and $\tau(A^2)\ll_\eps A^\eps$, replacing the denominator costs $O(1)$.  Apply the two halves of Lemma~\ref{lem:exact-R} in \eqref{eq:tau-harmonic}; the subtracted $1$ in the lower bound costs only $O(\log Q)$.
\end{proof}

A mass-only Rankin sieve with auxiliary scale $X=\e^t$, class mass $O(t^B)$, and product-tail budget $t^{B+1}\ll\log N$ can save at most $O((\log N)^{B/(B+1)})$.  Thus the intrinsic supply has the formal ceiling $3/4$.  This is only a ceiling inside that mass/product-tail model; it is not a nonexistence theorem for other uses of the class geometry.

\subsection{Shifted-divisor form}
\begin{lemma}[exact reorganization]\label{lem:shifted}
For $A=(M+1)/4$,
\[
 D\mid A^2\Longleftrightarrow R_0(D)\mid A
 \Longleftrightarrow M\equiv-1\pmod{4R_0(D)}.
\]
Consequently $n\notin\Av_X$ if and only if there are $D,M\geq1$ such that
\begin{equation}\label{eq:shifted}
 M\mid n+4D,\qquad M\leq X,\qquad
 M\equiv-1\pmod{4R_0(D)}.
\end{equation}
If $R=\prod p^{a_p}$, then
\begin{equation}\label{eq:fibers}
 R_0(D)=R\Longleftrightarrow D=R^2/s
 \quad\text{for some }s\mid\rad(R).
\end{equation}
\end{lemma}
\begin{proof}
At a prime, $D\mid A^2$ says $e_p(D)\leq2v_p(A)$, exactly the assertion $\lceil e_p(D)/2\rceil\leq v_p(A)$.  This proves the first line.  Lemma~\ref{lem:exact-R} changes $n\equiv-4D\pmod M$ into \eqref{eq:shifted}, in both directions.  Finally, if $v_p(R)=a$, the condition $\lceil v_p(D)/2\rceil=a$ allows $v_p(D)=2a$ or $2a-1$, independently at each $p\mid R$, which is \eqref{eq:fibers}.
\end{proof}
The congruence in \eqref{eq:shifted} implies $R_0(D)\leq(X+1)/4$ and hence $D\leq(X+1)^2/16$; no $D$-tail is omitted.

\section{Squares and the prime slice}\label{sec:prime}

\subsection{The square-class phenomenon}
\begin{lemma}[quadratic sign]\label{lem:jacobi}
If $M=4A-1$ and $D\mid A^2$, then
\[
 \left(\frac DM\right)=1,
 \qquad
 \left(\frac{-4D}M\right)=-1.
\]
Thus every unit class in $\Rcal(M)$ has Jacobi sign $-1$.  Every $n$ with $(n/M)\in\{0,1\}$ avoids $\Rcal(M)$; in particular, every perfect square avoids the complete intrinsic system for every $X$.
\end{lemma}
\begin{proof}
For an odd prime $p\mid A$, quadratic reciprocity and $M\equiv-1\pmod p$, $M\equiv3\pmod4$ give
\[
 \left(\frac pM\right)=\left(\frac Mp\right)
 (-1)^{(p-1)(M-1)/4}
 =\left(\frac{-1}p\right)(-1)^{(p-1)/2}=1.
\]
If $2\mid A$, then $M\equiv7\pmod8$ and $(2/M)=1$.  Multiplicativity gives $(D/M)=1$.  Since $(-1/M)=-1$ and $(4/M)=1$, the second identity follows.  Also $(A,M)=1$, so all intrinsic classes are units.
\end{proof}
In particular $|\Av_X(N)|\geq\lfloor\sqrt N\rfloor$.  Unit squares modulo the odd period
$P_X=\lcm\{M\leq X:M\equiv3\pmod4\}$ also give
\[
 \delta_X\geq\frac{\varphi(P_X)}{2^{\omega(P_X)}P_X}
 \geq\exp\{-O(X/\log X)\}.
\]
These bounds are too small to test a cubic exponent when $\log N\asymp L^4$.

\subsection{Harvesting prime classes}
We record the fixed-multiplier harvesting input in a self-contained form.
\begin{lemma}[prime harvest]\label{lem:prime-harvest}
For primes $\ell\equiv3\pmod4$, there are subsets $\mathcal C_\ell\subseteq\Rcal(\ell)$ such that
\begin{equation}\label{eq:harvest}
 \sum_{X^{1/2}<\ell\leq X}\frac{|\mathcal C_\ell|}{\ell}\asymp L^2.
\end{equation}
\end{lemma}
\begin{proof}
In a dyadic interval $x<\ell\leq2x$, put $z=x^{1/6}$ and take coprime $u,v\leq z$ with $\omega(uv)\leq D\log\log X$.  The progression
$\ell\equiv-1\pmod{4uv}$ makes $uv\mid(\ell+1)/4$, so Lemma~\ref{lem:multiplier} supplies the residue $-uv^{-1}\pmod\ell$.  Distinct reduced pairs give distinct residues: a collision implies
$\ell\mid uv'-u'v$, while $|uv'-u'v|<\ell^{2/3}<\ell$, and hence $u/v=u'/v'$.

The elementary harmonic estimates
\[
 \sum_{\substack{u,v\leq z\\(u,v)=1}}\frac1{uv}\asymp(\log z)^2,
 \qquad
 \sum_{u,v\leq z}\frac1{\varphi(u)\varphi(v)}\ll(\log z)^2
\]
follow respectively by M\"obius inversion and from $\sum_{n\leq z}1/\varphi(n)\ll\log z$.  Moreover, for fixed $1<t<2$,
\[
 \sum_{u,v\leq z}\frac{t^{\omega(u)+\omega(v)}}{uv}
 \ll_t(\log z)^{2t}.
\]
Rankin's inequality therefore shows that a fixed $D$ leaves a positive proportion of the first harmonic sum.  A modulus $4uv$ is then generated at most $2^{\omega(uv)}\leq(\log X)^{D\log2}$ times.  Bombieri--Vinogradov applies since $4uv\leq4x^{1/3}$ and absorbs this fixed logarithmic multiplicity.  It gives $\gg x\log x$ retained prime incidences in the block; division by $\ell\asymp x$ gives $\gg\log x$.  There are $\asymp L$ dyadic blocks between $X^{1/2}$ and $X$, proving the lower bound in \eqref{eq:harvest}.  Brun--Titchmarsh and the second harmonic estimate give $O(x\log x)$ incidences per block and hence the matching upper bound.
\end{proof}

\begin{lemma}[full prime mass]\label{lem:prime-mass}
There are $c,C>0$ such that
\begin{equation}\label{eq:prime-mass}
 cL^2\leq\sum_{\substack{\ell\leq X\\\ell\equiv3\ (4)}}\frac{F(\ell)}\ell\leq CL^2,
\end{equation}
where $\ell$ is prime.
\end{lemma}
\begin{proof}
The lower bound follows from Lemma~\ref{lem:prime-harvest}.  For the upper bound put $A=(\ell+1)/4$.  Lemma~\ref{lem:exact-R} and $2e+1\leq(e+1)(e+2)/2$ give
$F(\ell)\leq\tau(A^2)\leq\tau_3(A)$.  On $A\asymp Y$, write $A=abc$ and designate a largest coordinate as $c$, at a factor-three cost; then $ab\ll Y^{2/3}$.  For fixed $a,b$, Brun--Titchmarsh for the prime $4abc-1$ modulo $4ab$ gives
\[
 \#\{c:A\asymp Y,\ 4abc-1\text{ prime}\}
 \ll\frac{Y}{\varphi(4ab)\log Y}.
\]
Since $\varphi(4ab)\geq\varphi(a)\varphi(b)$ and
$\sum_{m\leq y}1/\varphi(m)\ll\log(2y)$,
\[
 \sum_{\substack{A\asymp Y\\4A-1\ \mathrm{prime}}}\tau(A^2)\ll Y\log Y.
\]
After division by $\ell\asymp Y$, each dyadic block costs $O(\log Y)$, and summing proves the upper bound.
\end{proof}

\begin{theorem}[two-sided prime-slice avoidance]\label{thm:prime-two-sided}
There are $c_1,C_2,C_3>0$ such that, for all large $X$ and $\log N\geq C_3L^3$,
\begin{equation}\label{eq:prime-two-sided}
 N\e^{-C_2L^2}\leq|\Av_X^{\rm prime}(N)|
 \leq N\e^{-c_1L^2}.
\end{equation}
Its natural density is $\e^{-\Theta(L^2)}$.
\end{theorem}
\begin{proof}
Chinese remaindering over distinct prime moduli gives exact density
\[
 V_X=\prod_{\substack{\ell\leq X\\\ell\equiv3\ (4)}}
 \left(1-\frac{F(\ell)}\ell\right).
\]
Lemma~\ref{lem:jacobi} gives $F(\ell)/\ell<1/2$.  Lemma~\ref{lem:prime-mass} and
$-2u\leq\log(1-u)\leq-u$ imply $V_X=\e^{-\Theta(L^2)}$.

For the finite interval put $a_\ell=F(\ell)/\ell$ and $\mu=\sum a_\ell=\Theta(L^2)$.  Odd and even Bonferroni truncations at degrees $r=O(\mu)$ bracket the avoider indicator.  Their omitted CRT main tails are bounded by
$\sum_{j>r}\mu^j/j!$; taking the constant in $r$ large makes each tail at most $V_X/4$.  An intersection indexed by a product $d$ of $j$ primes occupies $\prod_{\ell\mid d}F(\ell)$ classes modulo $d$, so its finite count is
\[
 N\prod_{\ell\mid d}a_\ell+O\!\left(\prod_{\ell\mid d}F(\ell)\right).
\]
The preceding proof of Lemma~\ref{lem:prime-mass}, before division by $\ell$, gives $W:=\sum_{\ell\leq X}F(\ell)\ll XL$.  Hence all rounding errors through degree $r$ total at most
$\sum_{j\leq r}W^j/j!=\exp\{O(L^3)\}$.  If $C_3$ is large this is at most $NV_X/4$.  The two Bonferroni bounds now give \eqref{eq:prime-two-sided}.
\end{proof}

Since the full avoiders are contained in the prime-slice avoiders, we obtain the promised unconditional full-system upper bound.
\begin{corollary}\label{cor:full-upper}
If $\log N\geq CL^3$, then
\[
 |\Av_X(N)|\ll N\exp\{-cL^2\}.
\]
The same upper bound holds for $\delta_X$ without a condition on $N$.
\end{corollary}

\subsection{A growing multi-shift certificate}
For $Y\geq1$, let $\Av_X^{(R\leq Y)}$ denote avoidance of every pair in \eqref{eq:shifted} with $R_0(D)\leq Y$.
\begin{theorem}[multi-shift prime certificate]\label{thm:multi-shift}
Fix $\eps>0$.  For all sufficiently large $X$, uniformly for $1\leq Y\leq L^{3-\eps}$ and $\log N\geq L^4$,
\[
 |\Av_X^{(R\leq Y)}(N)|
 \geq N\exp\{-C_\eps Y\log(2Y)\log L\}.
\]
The corresponding natural density has the same lower bound.  At $Y=L^{3-\eps}$ the exponent is $o(L^3)$.
\end{theorem}
\begin{proof}
Let
\[
 \mathcal D_Y=\{(R,D):R\leq Y,\ D=R^2/s,\ s\mid\rad(R)\},
 \qquad H=|\mathcal D_Y|.
\]
The standard Euler-product estimate
\[
 \sum_{R\leq Y}2^{\omega(R)}\ll Y\log(2Y)
\]
gives $H\ll Y\log(2Y)$.  For each prime $p\leq X$, $p\equiv3\pmod4$, forbid
\[
 B_p=\{-4D\pmod p:(R,D)\in\mathcal D_Y,\ p\nmid R\}.
\]
Every $B_p$ is nonzero and has size $g(p)\leq\min(H,p-1)$.  Avoiding all $B_p$ certifies the desired avoidance: a hit in \eqref{eq:shifted} has $(M,R)=1$, and some prime $p\equiv3\pmod4$ divides $M$ to odd exponent; then $p\leq X$, $p\nmid R$, and $n\equiv-4D\pmod p$, a contradiction.

For $p\leq2H$, choose the allowed residue $n\equiv0\pmod p$, and let $Q_0$ be their product.  Then $\log Q_0=O(H)$.  In $n=Q_0m$, every remaining forbidden set has $g(p)\leq H<p/2$, and
\[
 \mu_Y=\sum_{2H<p\leq X}\frac{g(p)}p\ll H\log L.
\]
Odd Bonferroni at degree $r\asymp\mu_Y$ retains a fixed fraction of the Euler product, at least $\exp(-2\mu_Y)$.  Since $W_Y=\sum g(p)\leq H\pi(X)$, the total rounding error is
\[
 \exp\{O(r\log W_Y)\}=\exp\{O(HL\log L)\}.
\]
For $Y\leq L^{3-\eps}$ this is $\exp\{o(L^4)\}$ and is negligible under $\log N\geq L^4$.  Restoring $Q_0$ proves the claimed bound.  On complete periods there is no rounding error, proving the density statement.
\end{proof}

\section{The cubic-mass hypothesis and its conditional consequence}\label{sec:hpf}
Put
\[
 \mu_X=\sum_{\substack{M\leq X\\M\equiv3\ (4)}}\frac{F(M)}M\asymp L^3.
\]
The following is the literal hypothesis proposed for partition-free assembly.  Its all-range quantifier in $N$ is essential.

\begin{hypothesis}[the literal cubic-mass majorant $\PF$]\label{hyp:PF}
There are absolute $c,C>0$ such that, for every sufficiently large $X$, every $J\geq C\mu_X$, and every $N$ satisfying
\[
 J\log X\leq\tfrac12\log N,
\]
there is a nonnegative arithmetic function $\nu$ on $[1,N]$ which is at least $1$ on every member of $\Av_X(N)$ and satisfies
\[
 \frac1N\sum_{n\leq N}\nu(n)\ll\exp(-c\mu_X).
\]
Its evaluation uses exact congruence counts modulo lcms of at most $J$ participating moduli; every such lcm is at most $X^J\leq N^{1/2}$.  The coefficient sum and number of congruence-count terms are at most $\exp(O(J))$, so aggregate rounding errors are required to be negligible under the displayed degree budget.
\end{hypothesis}
This hypothesis concerns the complete system, not arbitrary subfamilies.  The latter formulation is already false: all $D=1$ classes with $3\mid M$ lie in $n\equiv2\pmod3$, while the complement has density $2/3$ despite growing class mass.

\begin{theorem}[formal conditional $3/4$ bound]\label{thm:conditional}
If $\PF$ holds, then for some $c>0$ the number $E_{\rm all}(N)$ of exceptional denominators $n\leq N$ satisfies
\[
 E_{\rm all}(N)\leq N\exp\{-c(\log N)^{3/4}\}.
\]
\end{theorem}
\begin{proof}
Put $t=\alpha(\log N)^{1/4}$, $X=\e^t$.  By Theorem~\ref{thm:cubic}, $\mu_X\asymp t^3$.  Take $J=C\mu_X$.  Then
$J\log X\ll\alpha^4\log N$, so small fixed $\alpha$ meets the degree condition.  Hypothesis~\ref{hyp:PF} bounds the avoiders by $N\exp(-c\mu_X)$.  Every exceptional denominator is an avoider, because every intrinsic class is forced soluble by Lemma~\ref{lem:multiplier}.
\end{proof}
The implication is valid, but its antecedent will now be disproved.

\section{Rough moduli and refutation of \texorpdfstring{$\PF$}{H-PF}}\label{sec:refutation}
For primes $q$ define
\[
 b_q=\frac{L^3}{q\log z},\qquad
 W_{\boldsymbol a}(M)=\1_{P^-(M)>z}
 \exp\!\left(2\sum_{p\mid M}a_p\right).
\]
The sum in the exponential is over distinct prime factors.  An \emph{atom} is a pair $(M,r)$ with $M\leq X$, $M\equiv3\pmod4$, and $r\in\Rcal(M)$; it is the event $n\equiv r\pmod M$.  Thus a modulus $M$ contributes exactly $F(M)$ atoms.

\begin{theorem}[weighted rough-modulus mean]\label{thm:rough-mean}
Uniformly for $3\leq z\leq X^{1/10}$ and charges $0\leq a_p\leq1$, put
$E=\sum_{z<p\leq X}a_p/p$.  Then
\begin{equation}\label{eq:rough-mean}
 \sum_{\substack{M\leq X\\M\equiv3\ (4)}}
 \frac{F(M)}M W_{\boldsymbol a}(M)
 \ll \e^{CE}\frac{L^3}{\log z}.
\end{equation}
In particular, if $a_p=O(L^3/(p\log z))$ and
$z=L^3/\sqrt{\log L}$, then $E=o(1)$ and the left side is $o(L^3)$.
\end{theorem}
\begin{proof}
Write $M=4m-1$ and define the multiplicative function
\[
 w(p^\nu)=\begin{cases}0,&p\leq z,\\ \e^{2a_p},&p>z,
 \end{cases}\qquad w(1)=1.
\]
Lemma~\ref{lem:exact-R} gives
$F(4m-1)W_{\boldsymbol a}(4m-1)\leq w(4m-1)\tau_3(m)$.
Apply the Nair--Tenenbaum estimate in Henriot's discriminant-uniform form \cite{NairTenenbaum,Henriot} to the determinant-one linear forms
$Q_1(t)=4t-1$, $Q_2(t)=t$ and the multiplicative function
$\mathcal F(u,v)=w(u)\tau_3(v)$.  The class parameters are uniform in the charges.  On a dyadic interval of length $V$, the local factors give
\[
 V\prod_{p\ll V}(1-2/p)(1+w(p)/p)(1+3/p)
 \ll V\frac{\log^2(2V)}{\log z}
 \exp\!\left(C\sum_{z<p\ll V}\frac{a_p}p\right).
\]
Indeed the coefficient of $1/p$ in the logarithm is $1$ for $p\leq z$ and $2+O(a_p)$ for $p>z$.  The first possible rough value exceeds $z$, so the lowest nonempty blocks obey the same estimate after changing the constant.  Dyadic summation and partial summation yield
\[
 \sum\frac{F(M)W_{\boldsymbol a}(M)}M
 \ll\frac{\e^{CE}}{\log z}
 \left(L^2+\int_z^X\frac{(\log t)^2}{t}\,dt\right),
\]
which proves \eqref{eq:rough-mean}.  Finally
$E\ll L^3/(z\log z)=o(1)$ for the stated charges and cutoff.
\end{proof}

The pointwise neighborhood estimate needs a small maximal-order correction.  Fix a sufficiently large absolute $C_d$ and put
\[
 u_q=\frac{\log(2q)}q
 \exp\!\left(\frac{C_d\log(2q)}{\log\log(3q)}\right).
\]
For an atom modulus divisible by $q$, set
\[
 T_q(z;\boldsymbol a)=
 \sum_{\substack{M\leq X,\ M\equiv3\ (4)\\P^-(M)>z,\ q\mid M}}
 \frac{F(M)}M\exp\!\left(2\sum_{p\mid M}a_p\right).
\]

\begin{theorem}[softened neighborhood bound]\label{thm:soft-neighbor}
Let
\begin{equation}\label{eq:zchoice}
 z=\frac{L^3\log\log L}{\log L}.
\end{equation}
For all large $X$ there is an absolute $A>0$ such that
$a_q=A(b_q+u_q)$ for primes $z<q\leq X$ satisfies
\[
 T_q(z;\boldsymbol a)\leq a_q,
 \quad \max_{q>z}a_q=o(1),
 \quad\sum_{q>z}\frac{a_q}q=o(1),
\]
and
\begin{equation}\label{eq:weighted-o}
 \sum_{\substack{M\leq X\\M\equiv3\ (4)}}
 \frac{F(M)}M W_{\boldsymbol a}(M)
 \ll\frac{L^3}{\log z}=o(L^3).
\end{equation}
The same conclusions hold for $z=L^3g(L)/\log L$ if
$g(L)\to\infty$ and $g(L)=o(\log L)$.
\end{theorem}
\begin{proof}
Use the maximal-order bound
\begin{equation}\label{eq:maxdiv}
 \tau_3(n)\leq\exp\!\left(\frac{C_d'\log(2n)}{\log\log(3n)}\right).
\end{equation}
Fix $q>z$, write $M=qk$, and choose $r\in\{1,3\}$ with $qr\equiv3\pmod4$.  With $k=4t+r$,
\[
 \frac{qk+1}{4}=qt+c,\qquad c=\frac{qr+1}{4},\qquad4c-qr=1.
\]
The cofactor and shifted factor are determinant-one linear forms, and the coefficient norm of their product is at most $13q$.  Also
$W_{\boldsymbol a}(qk)\leq\e^{2a_q}w(k)$ and
$F(qk)\leq\tau_3((qk+1)/4)$.

Split at $k=C_0q^{1/2}$, where $C_0$ is large enough for Henriot's coefficient-size condition.  In the lower range, \eqref{eq:maxdiv}, roughness, and $\sum_{k\leq y}1/k\ll\log(2y)$ give
\[
 \frac1q\sum_{k<C_0q^{1/2}}
 \frac{F(qk)W_{\boldsymbol a}(qk)}k\ll u_q.
\]
Here $a_p=O(1/\log\log L+p^{-1+o(1)})$, and a rough $qk$ has at most $\log(qk)/\log z$ prime factors, so its exponential weight is absorbed by enlarging the constant in \eqref{eq:maxdiv}.  This includes $k=1$.

Decompose the high range dyadically in $t$.  Henriot's theorem applies to the two determinant-one forms and $w(4t+r)\tau_3(qt+c)$, giving
\[
 \sum_{V<k\leq2V}w(k)\tau_3((qk+1)/4)
 \ll V\frac{(\log(2V))^2}{\log z}.
\]
After division by $qk$ and summation over blocks, this is
$O(L^3/(q\log z))=O(b_q)$.  Thus $T_q\ll b_q+u_q$ uniformly once
$\max a_p\leq1$ and $\sum a_p/p\leq1$.  Choose the fixed $A$ larger than this implicit constant; the estimates below verify those two conditions for all sufficiently large $X$, closing the argument without circularity.

Indeed
\[
 \sum_{q>z}\frac{b_q}q\ll\frac{L^3}{z\log z}=o(1).
\]
For large $q$, the exponential in $u_q$ is at most $q^{1/4}$, whence
\[
 \sum_{q>z}\frac{u_q}{q}\ll
 \sum_{n>z}\frac{\log(2n)}{n^{7/4}}=o(1).
\]
The same estimates give $\max a_q=o(1)$.  Theorem~\ref{thm:rough-mean} proves \eqref{eq:weighted-o}.  Replacing $\log\log L$ by $g(L)$ changes the first displayed bound to $O(1/g(L))$ and nothing else.
\end{proof}

\begin{theorem}[natural-density lower bound]\label{thm:density-lower}
With \eqref{eq:zchoice},
\begin{equation}\label{eq:density-lower}
 \delta_X\geq
 \exp\!\left\{-O\!\left(\frac{L^3\log\log L}{\log L}\right)\right\}
 =\exp\{-o(L^3)\}.
\end{equation}
More generally, the last cutoff in Theorem~\ref{thm:soft-neighbor} gives
$\delta_X\geq\exp\{-O(L^3g(L)/\log L)\}$.
\end{theorem}
\begin{proof}
First impose $n\equiv0\pmod p$ for every odd prime $p\leq z$.  Work on the finite uniform CRT space modulo the lcm of their product and all participating moduli.  The conditioning costs
$\exp\{-\vartheta(z)+O(1)\}=\exp\{-O(z)\}$.  If an intrinsic atom has a modulus divisible by such a $p$, it is now impossible: its residue $-4D$ is a unit modulo $M$, since $(D,M)=1$.  Every surviving distinct atom has a $z$-rough modulus $M$, remains uniform, and has probability exactly $1/M$.

Join two atoms when their moduli share a prime and set
\[
 x_A=\frac1M\exp\!\left(2\sum_{q\mid M}a_q\right)
\]
for an atom of modulus $M$.  Roughness and $\max a_q=o(1)$ give $x_A<1/2$.  Theorem~\ref{thm:soft-neighbor} bounds the total activity of the neighbors of $A$ by $\sum_{q\mid M}a_q$.  Therefore
\[
 \prod_{B\sim A}(1-x_B)
 \geq\exp\!\left(-2\sum_{B\sim A}x_B\right)
 \geq\exp\!\left(-2\sum_{q\mid M}a_q\right),
\]
so $\Pr(A)\leq x_A\prod_{B\sim A}(1-x_B)$.  The quantitative local lemma gives
\[
 \begin{aligned}
 \Pr(\text{no surviving atom}\mid\text{quarantine})
 &\geq\prod_A(1-x_A)\\
 &\geq\exp\!\left(-2\sum_Ax_A\right)
 \geq\exp\{-O(L^3/\log z)\}.
 \end{aligned}
\]
where the last inequality is \eqref{eq:weighted-o}.  Multiplying by the quarantine cost proves \eqref{eq:density-lower}.
\end{proof}

\begin{corollary}[refutation of the literal hypothesis]\label{cor:refute}
Hypothesis $\PF$ as stated in Hypothesis~\ref{hyp:PF} is false.
\end{corollary}
\begin{proof}
Theorem~\ref{thm:cubic} gives $\mu_X\geq c_0L^3$.  If $\PF$ held with constants $c,C$, choose $J\geq C\mu_X$.  For each fixed sufficiently large $X$, take arbitrarily large multiples $N=mP_X$ of the full period satisfying
$J\log X\leq\tfrac12\log N$.  Periodicity gives
$|\Av_X(N)|/N=\delta_X$.  Since the majorant is at least one on every avoider, $\PF$ would imply
\[
 \delta_X\ll\e^{-c\mu_X}\leq\e^{-cc_0L^3}.
\]
This contradicts Theorem~\ref{thm:density-lower} for all sufficiently large $X$.
\end{proof}
The scope is exact: this refutes the all-range quantifiers of $\PF$, not the conclusion of Theorem~\ref{thm:conditional} by every possible route and not the critical-window statement below.

\section{The critical window and transfer walls}\label{sec:window}
\begin{hypothesis}[critical-window variant $\PFp$; open]\label{hyp:PFprime}
There are $c,C>0$ and $0<c_-<c_+$, with $c_-$ large enough for the degree budget, such that the majorant in Hypothesis~\ref{hyp:PF} exists whenever
\[
 c_-L^4\leq\log N\leq c_+L^4,
 \qquad J\geq C\mu_X,
 \qquad J\log X\leq\tfrac12\log N.
\]
No assertion is made outside this window.
\end{hypothesis}
This remains open and would still imply Theorem~\ref{thm:conditional}.  Corollary~\ref{cor:refute} uses arbitrarily large complete-period multiples at fixed $X$, which need not lie in this window.

\subsection{Cylinder and tensor obstructions}
\begin{theorem}[all-avoider cylinder wall]\label{thm:cylinder}
If a nonempty union $\mathcal U$ of complete classes modulo $Q$ is contained in the full intrinsic avoider set, then
\[
 \prod_{\substack{\ell\leq X,\ \ell\equiv3\ (4)\\\ell\ \mathrm{prime}}}\ell\mid Q,
 \qquad \log Q\geq(1/2+o(1))X.
\]
\end{theorem}
\begin{proof}
Choose one class $a\pmod Q$ in $\mathcal U$.  For each prime $\ell\equiv3\pmod4$, the divisor $D=1$ gives the forbidden class $-4\pmod\ell$.  If $\ell\nmid Q$, the Chinese remainder theorem produces an integer lying both in $a\pmod Q$ and in that forbidden class, a contradiction.  Thus every such prime divides $Q$, and the prime number theorem in the progression $3\pmod4$ gives the estimate.
\end{proof}

The CRT space is a product over prime powers, but the avoider indicator is not a tensor product: its atomic factors overlap.
\begin{lemma}[explicit non-tensor obstruction]\label{lem:nontensor}
Already at $X=15$, the complete avoider indicator modulo $1155=3\cdot5\cdot7\cdot11$ does not factor into functions of the four prime coordinates.
\end{lemma}
\begin{proof}
The residues
\[
\begin{array}{c|cc|cc|c}
 n&n\bmod3&n\bmod5&n\bmod7&n\bmod11&\1_{\Av_{15}}(n)\\ \hline
231&0&1&0&0&1\\
616&1&1&0&0&1\\
693&0&3&0&0&1\\
1078&1&3&0&0&0
\end{array}
\]
form a two-coordinate rectangle with the other coordinates fixed.  A tensor product containing the first three corners contains the fourth.  To check the final column without an enumeration, Lemma~\ref{lem:exact-R} gives
\[
 \Rcal(3)=\{2\},\quad \Rcal(7)=\{3,5,6\},\quad
 \Rcal(11)=\{7,8,10\},\quad \Rcal(15)=\{7,11,13,14\}.
\]
The first three rows miss these four sets.  The last row is $13\pmod{15}$ and is killed by $D=8\mid4^2$.
\end{proof}
Thus Fourier transformation produces a hypergraph convolution, not a product of local Fourier transforms.

\subsection{The corrected finite-transfer criterion}
After the odd-prime quarantine, write $n=P_zm$.  For $U\subseteq\mathbb Z/d\mathbb Z$, put $1_{U,d}(m)=\1_{m\bmod d\in U}$ and $\rho(U)=|U|$.
\begin{lemma}[residue-set transfer]\label{lem:residue-transfer}
Let $\mathcal F_X$ be a specified family of surviving atoms and $I_X$ a finite index set.  Suppose
\[
 B_X(m)=\sum_{\alpha\in I_X}c_\alpha1_{U_\alpha,d_\alpha}(m)
 \leq\1_{\{\text{no atom of }\mathcal F_X\text{ at }m\}}
\]
pointwise, where each $d_\alpha$ divides the surviving CRT period and $(d_\alpha,P_z)=1$.  If
\[
 \begin{gathered}
 \mathbb E_{\rm CRT}B_X\geq\e^{-o(L^3)},\qquad
 \log\max_\alpha d_\alpha=o(L^4),\\
 \log\mathcal R_X=o(L^4),\qquad
 \mathcal R_X=\sum_\alpha|c_\alpha|\rho(U_\alpha).
 \end{gathered}
\]
then, uniformly for fixed $0<c_-<c_+$ and
$c_-L^4\leq\log N\leq c_+L^4$,
\[
 |\Av(\mathcal F_X;N)|\geq N\e^{-o(L^3)}.
\]
\end{lemma}
\begin{proof}
Put $H=\lfloor N/P_z\rfloor$.  A set of $\rho(U_\alpha)$ classes modulo $d_\alpha$ has count
\[
 H\frac{\rho(U_\alpha)}{d_\alpha}+O(\rho(U_\alpha)).
\]
Hence
$\sum_{m\leq H}B_X(m)=H\mathbb E_{\rm CRT}B_X+O(\mathcal R_X)$.
The main term has logarithm $\log N-O(z)-o(L^3)=\log N-o(L^3)$, while the error has logarithm $o(L^4)$ and is negligible because $\log N\geq c_-L^4$.  The pointwise inequality and $P_z=\exp\{O(z)\}=\exp\{o(L^3)\}$ finish the proof.
\end{proof}
For an atom monomial, compatibility gives one residue class and incompatibility gives none, so $\rho\leq1$.  The following is the isolated low-degree criterion; the factor $\rho(U)$ in the general ledger is essential.
\begin{corollary}[monomial finite-transfer criterion]\label{cor:monomial-transfer}
Suppose real coefficients $c_S$, supported on sets of at most $r=o(L^3)$ surviving atoms, make
\[
 B_X(n)=\sum_Sc_S\prod_{A\in S}\1_A(n)
 \leq\1_{\{\text{no surviving atom at }n\}}
\]
pointwise.  If its exact CRT mean is at least $\e^{-o(L^3)}$ and
$\log\sum_S|c_S|=o(L^4)$, then
$|\Av_X(N)|\geq N\e^{-o(L^3)}$ uniformly for $\log N\asymp L^4$.
\end{corollary}
\begin{proof}
Apply Lemma~\ref{lem:residue-transfer}.  Compatible monomials are single classes modulo lcms $d_S\leq X^r=\exp\{o(L^4)\}$, incompatible monomials vanish, and the weighted ledger is at most $\sum_S|c_S|$.
\end{proof}

\subsection{The prime side transfers}
Let $S_\ell=P_z^{-1}\Rcal(\ell)\pmod\ell$.  Choose $h(L)\to\infty$ slowly and put $Y=L^4/h(L)$, so $z<Y<X$.  Define
\[
 d_0=\prod_{\substack{z<\ell\leq Y\\\ell\equiv3\ (4)}}\ell,
 \qquad
 U_0=\prod_{\substack{z<\ell\leq Y\\\ell\equiv3\ (4)}}
 ((\mathbb Z/\ell\mathbb Z)\setminus S_\ell),
 \qquad T_0=1_{U_0,d_0}.
\]
Here and below $\ell$ is prime.  Then $\log d_0=O(Y)=o(L^4)$.  Let $H_P$ count violated prime atoms with $Y<\ell\leq X$, put $\mu_P=\mathbb EH_P=\Theta(L^2)$, let $r$ be the least odd integer at least $8\mu_P$, and define
\[
 Q_r(H_P)=\sum_{j=0}^r(-1)^j\binom{H_P}{j},\qquad B_P=T_0Q_r(H_P).
\]

\begin{theorem}[prime-side two-level minorant]\label{thm:prime-transfer}
The function $B_P$ is pointwise at most the surviving prime-system avoider indicator.  Its tail degree is $O(L^2)$, every term has modulus logarithm $o(L^4)$,
\[
 \mathbb E_{\rm CRT}B_P\geq\e^{-O(L^2)},
 \qquad
 \log\sum_\alpha|c_\alpha|\rho(U_\alpha)=o(L^4).
\]
Thus Lemma~\ref{lem:residue-transfer} transfers a lower bound $N\e^{-o(L^3)}$ in the critical window.  The sharper $N\e^{-O(L^2)}$ bound is already Theorem~\ref{thm:prime-two-sided}.
\end{theorem}
\begin{proof}
If $T_0=0$ the pointwise assertion is immediate; if $T_0=1$, it is odd Bonferroni.  For $h\geq1$,
\[
 \sum_{j=0}^r(-1)^j\binom{h}{j}=-\binom{h-1}{r}.
\]
Distinct prime coordinates are independent and atoms at one coordinate are disjoint, so the loss from the exact tail void is at most
\[
 \mathbb E\binom{H_P}{r+1}\leq\frac{\mu_P^{r+1}}{(r+1)!}
 \leq\left(\frac{\e\mu_P}{r+1}\right)^{r+1}.
\]
The tail Euler product is at least $\e^{-2\mu_P}$; the last display is a small fraction of it.  This proves the mean bound, including the independent low tensor.

Expand only the tail factor.  If $W_P=\sum_{Y<\ell\leq X}F(\ell)$, the maximal-order divisor bound gives $W_P=\exp\{O(L)\}$.  Each term has at most $\rho(U_0)$ residue classes, so
\[
 \mathcal R_X\leq\rho(U_0)\sum_{j\leq r}\frac{W_P^j}{j!},
 \qquad \log\mathcal R_X\leq o(L^4)+O(rL)=o(L^4).
\]
Its modulus divides $d_0X^r$, giving the other budget.  Apply Lemma~\ref{lem:residue-transfer}; the quarantine costs $\e^{-O(z)}=\e^{-o(L^3)}$.
\end{proof}

\subsection{The composite frontier}
Delete a composite atom if one of its prime projections is already a prime intrinsic atom; retain the residual family $\mathcal C^\circ$.  Condition on $T_0=1$, and let $H$ count the tail prime atoms together with $\mathcal C^\circ$.  Put
\[
 \Lambda=\frac{L^3}{\log L}.
\]
For $A\in\mathcal C^\circ$, conditioning on $T_0=1$ gives
\[
 \Pr(A\mid T_0=1)=\frac1{M_A}
 \prod_{\substack{p\mid M_A\\z<p\leq Y}}
 \frac{p}{p-F(p)}=\frac{1+o(1)}{M_A}
\]
uniformly: $F(p)=p^{o(1)}$, $p>z=L^{3+o(1)}$, and
$\omega(M_A)\leq L/\log z$.  The rough mean therefore gives
$\mathbb EH=O(\Lambda)$.  The conditional full void after the quarantine is at least $\e^{-O(\Lambda)}$ by the proof of Theorem~\ref{thm:density-lower}; dividing by $\Pr(T_0=1)\leq1$ gives
$\Pr(H=0\mid T_0=1)\geq\e^{-O(\Lambda)}$.

\begin{proposition}[factorial-moment reduction]\label{prop:factorial}
Suppose fixed $C,D>0$, with $D$ sufficiently large in terms of $C$ and of the constant in the conditioned void lower bound (37.18) of the source notes, and an even $m\in[D\Lambda,D\Lambda+2]$ satisfy
\begin{equation}\label{eq:factorial-target}
 \mathbb E(H)_m\leq(C\Lambda)^m.
\end{equation}
Then a two-level pointwise minorant satisfies Lemma~\ref{lem:residue-transfer} and proves
$|\Av_X(N)|\geq N\e^{-o(L^3)}$ uniformly for $\log N\asymp L^4$.
\end{proposition}
\begin{proof}
Use $B=T_0Q_{m-1}(H)$.  Since $m-1$ is odd, Bonferroni makes it a pointwise minorant.  The preceding binomial identity and
$\binom{h-1}{m-1}\leq\binom{h}{m}$ give
\[
 \mathbb E(Q_{m-1}(H)\mid T_0=1)
 \geq\Pr(H=0\mid T_0=1)-\frac{(C\Lambda)^m}{m!}.
\]
By Stirling, the second term is at most $(\e C/D)^m$, a small fraction of the first for large $D$.  There are $\exp\{O(L)\}$ atoms.  Thus every term, after multiplication by $T_0$, is a residue set and
\[
 \begin{aligned}
 \deg B&=O(\Lambda)=o(L^3),\\
 \log(\mathrm{modulus})&\leq Y+O(\Lambda L)=o(L^4),\\
 \log(\mathrm{ledger})&\leq Y+O(\Lambda L)=o(L^4).
 \end{aligned}
\]
Lemma~\ref{lem:residue-transfer} applies.
\end{proof}

No proved estimate gives \eqref{eq:factorial-target}.  The exact first obstruction is residue dispersion.  For a high prime $p$ and $a\pmod p$, set
\[
 u_{p,a}=\sum_{\substack{A\in\mathcal C^\circ\\p\mid M_A,\ r_A\equiv a\ (p)}}
 \Pr(A\mid T_0=1),\qquad t_p=\sum_a u_{p,a}.
\]
For squarefree pairs whose gcd is $p$, their pair contribution is bounded by a sub-sum of
$p\sum_a u_{p,a}^2$.  Even granting the ideal incidence estimate $t_p\ll\Lambda/p$, the available $\ell^1$ information yields only
\begin{equation}\label{eq:logloss}
 \sum_{z<p\leq X}p\sum_a u_{p,a}^2
 \leq\sum_{z<p\leq X}pt_p^2
 \ll\Lambda^2\sum_{z<p\leq X}\frac1p
 \asymp\Lambda^2\log L.
\end{equation}
The live dispersion frontier is
\begin{equation}\label{eq:dispersion}
 \sum_{z<p\leq X}p\sum_a u_{p,a}^2=O(\Lambda^2),
\end{equation}
followed by higher-codegree analogues through $m\asymp\Lambda$.  Equation \eqref{eq:logloss} is not a counterexample to \eqref{eq:dispersion}; it identifies the factor lost by present information.

Scott--Sokal's cluster expansion \cite{ScottSokal} controls a numerical partition function, not the needed pointwise polynomial.  Indeed the literal independent-set polynomial in violated-event indicators fails pointwise: if the induced graph is a five-vertex path, its alternating independent-set sum is $1-5+6-1=1$ although the void indicator is zero.  Thus $\PFp$, the factorial-moment estimate, and the residue-resolved hierarchy beginning at \eqref{eq:dispersion} remain \textsc{open}.

\section{Provisional record context}\label{sec:records}
This section is deliberately segregated.  Both records are \textsc{claimed/provisional}; neither is used in the refutation of $\PF$.

\subsection{The provisional exceptional-set record}
Let $E(N)$ count exceptional prime denominators $p\leq N$.
\begin{theorem}[\textsc{claimed/provisional}]\label{thm:record}
There is an absolute $c>0$ such that
\[
 E(N)\ll N\exp\{-c(\log N)^{2/3}(\log\log N)^{1/3}\}.
\]
\end{theorem}
\begin{proof}[Internal argument; provisional status retained]
Put
\[
 \begin{aligned}
 \mathcal K(K)&=\{k\leq K:k\equiv1\pmod4\},\qquad
 L_K=\lcm_{k\in\mathcal K(K)}k,\\
 h(\mathcal J)&=\sum_{k\in\mathcal J}\frac{\varphi(k)}{k^2}.
 \end{aligned}
\]
The multiplier form of Lemma~\ref{lem:multiplier}, with $M=k\ell$, has the following class-mass consequence.  Uniformly for $K\leq(\log X)^A$, every subfamily $\mathcal J\subseteq\mathcal K(K)$ containing $1$, and every $c\pmod{24L_\mathcal J}$ with $(c,24L_\mathcal J)=1$, there are distinct forced classes modulo primes $X^{1/2}<\ell\leq X$ whose counts $f_c(\ell)$ satisfy
\begin{equation}\label{eq:mult-mass}
 \sum_{X^{1/2}<\ell\leq X}\frac{f_c(\ell)}\ell
 \asymp(\log X)^2h(\mathcal J).
\end{equation}
Here the classes arise from
\[
 H<u,v\leq\ell^{1/3},\quad (u,v)=(uv,k)=1,
 \quad uv\mid(k\ell+1)/4,
 \quad k\mid u+cv,
\]
with $H=K^{10}$ and $\omega(uv)\leq D\log\log X$.

For completeness, the harmonic lattice input behind \eqref{eq:mult-mass} is obtained by partitioning $(H,z]$ into fixed-ratio boxes.  M\"obius inversion of $(u,v)=1$ gives, in $I\times J$ with lengths $\eta U,\eta V$,
\[
 \begin{aligned}
 \#\{(u,v)\in I\times J:(u,v)=(uv,k)=1,\ k\mid u+cv\}
 &=\eta^2UV\frac{\varphi(k)}{k^2}
   \prod_{p\nmid k}(1-p^{-2})\\
 &\quad+O((U+V)\tau(k)\log(2UV)).
 \end{aligned}
\]
After division by $UV$ and summation over boxes, the boundary/main ratio is $O(K^2/H)$, uniformly in $k,c$.  Shiu's theorem \cite{Shiu}, applied successively in the two variables to $t^{\omega(n)}$ and $n/\varphi(n)$, gives
\[
 \sum^*\frac{t^{\omega(u)+\omega(v)}}{uv}
 \ll_t\frac{\varphi(k)}{k^2}(\log z)^2(\log z)^{2(t-1)},
 \qquad
 \sum^*\frac1{\varphi(u)\varphi(v)}
 \ll\frac{\varphi(k)}{k^2}(\log z)^2.
\]
Rankin's inequality removes the high-$\omega$ tail after fixed large $D$.  Summing over $k$ gives the factor $h(\mathcal J)$.

In a dyadic prime interval, the divisor condition is the progression
$\ell\equiv-k^{-1}\pmod{4uv}$.  Colliding classes imply
$\ell\mid uv'-u'v$ and hence equality of reduced pairs; different $k,k'$ would force the common $uv$ to divide $(k-k')/4$, impossible because $uv>H^2>K$.  A fixed modulus has multiplicity at most
$K2^{\omega(uv)}=(\log X)^{O(1)}$.  Bombieri--Vinogradov absorbs it and gives the lower half of \eqref{eq:mult-mass}; Brun--Titchmarsh and the second weighted lattice estimate give the upper half.  This proves \eqref{eq:mult-mass}.

Now choose $K=\lfloor\delta\log N\rfloor$ and partition primes by reduced classes modulo $M_0=24L_K$.  Since $\log M_0\leq(1+o(1))K$, each subsequence has length $Y\geq N^{1-2\delta}$.  Put
\[
 X=\exp\{\alpha(\log N/\log\log N)^{1/3}\}.
\]
For a fixed subsequence, the larger-sieve denominator is
\[
 S_c=\sum_{\substack{s\leq Y^{1/2}\\s\mid P}}\mu^2(s)
 \prod_{\ell\mid s}\frac{f_c(\ell)}{\ell-f_c(\ell)},
 \quad P=\prod_{X^{1/2}<\ell\leq X}\ell.
\]
Here $f_c(\ell)\leq K\ell^{2/3}<\ell/2$ for large $N$.  Writing $G_c$ for the untruncated Euler product, Rankin's trick with $v=1/\log X$ gives
\[
 \frac{G_c-S_c}{G_c}
 \leq\exp\{-\log Y/(2\log X)+C(\log X)^2\log K\}.
\]
Small $\alpha$ makes this at most $1/2$.  Since
$h(\mathcal K(K))\asymp\log K$, \eqref{eq:mult-mass} gives
$S_c\geq\tfrac12\exp\{a(\log X)^2\log K\}$.  Summing the uniform bound $O(Y/S_c)$ over the reduced subsequences proves the theorem.
\end{proof}

\begin{remark}[priority and verification status]\label{rem:priority}
Theorem~\ref{thm:record} remains \textup{CLAIMED}/\allowbreak\textup{PROVISIONAL} pending external checking and primary-literature verification.  Vaughan's primary paper \cite{Vaughan} was located by its confirmed DOI but its PDF was access-blocked.  Pomerance--Weingartner's 2025 reconstruction \cite[\S4]{PomeranceWeingartner} cites Vaughan as the state of the art, and the fresh search found no post-1970 exponent improvement.  The draft therefore does not claim a direct verification of Vaughan's argument or a settled record; the search result is not a literature guarantee.
\end{remark}

\subsection{The pruned cubic prime slice}
The current harvesting frontier uses the incidence
\[
 r_\mathcal J(u,v;c)=\#\{k\in\mathcal J:k\mid u+cv\}.
\]
\begin{lemma}[incidence second moment; provisional context]\label{lem:incidence}
For $K\geq K_0$, $H=K^{10}$, and $z>H^2$,
\[
 \sum_{\substack{H<u,v\leq z\\(u,v)=1}}
 \frac{r_\mathcal J(u,v;c)^2}{uv}
 \ll(\log z)^2(1+\log K)^3,
\]
uniformly for $\mathcal J\subseteq\mathcal K(K)$ and
$c\pmod{24L_\mathcal J}$ with $(c,24L_\mathcal J)=1$.
\end{lemma}
\begin{proof}
Expand the square.  For $k,k'\in\mathcal J$, both divisibilities are the congruence $[k,k']\mid u+cv$.  The preceding box calculation, now with $d=[k,k']\leq K^2$, gives
\[
 \sum^*\frac1{uv}\ll\frac{\varphi(d)}{d^2}(\log z)^2
       +\frac{\tau(d)(\log z)^2}{H}.
\]
Furthermore
\[
 \sum_{k,k'\leq K}\frac{\varphi([k,k'])}{[k,k']^2}
 \leq\sum_{k,k'\leq K}\frac{(k,k')}{kk'}
 \ll(1+\log K)^3,
\]
by writing $k=da,k'=db$.  Since $\tau([k,k'])\leq2K$, all boundary terms total $O(K^3(\log z)^2/H)$, proving the result.
\end{proof}

\begin{theorem}[unconditional pruned cubic prime slice; \textsc{claimed/provisional}]\label{thm:pruned}
Fix $0<\kappa<1/240$ and let $K=X^\kappa$.  For every
$\mathcal J\subseteq\mathcal K(K)$ containing $1$ and every
$c\pmod{24L_\mathcal J}$ with $(c,24L_\mathcal J)=1$, retain only triples in \eqref{eq:mult-mass} satisfying
$r_\mathcal J(u,v;c)\leq(\log X)^4$.  Then (i) the restriction of the
$H_{k\rm BV}$ error sum to these retained triples is
$o(x\log x\,h(\mathcal J))$ unconditionally, uniformly in all the
quantified data; and (ii) their distinct class counts
$f_c^{\rm good}(\ell)$ satisfy, uniformly in these data,
\[
 \sum_{X^{1/2}<\ell\leq X}\frac{f_c^{\rm good}(\ell)}\ell
 \asymp(\log X)^2h(\mathcal J).
\]
For the full $\mathcal J=\mathcal K(X^\kappa)$, this is $\asymp(\log X)^3$.
\end{theorem}
\begin{proof}[Internal argument; provisional status retained]
The harmonic lattice proof above remains uniform whenever $z>K^{20}$: its box error is $O(K^2/H)$, and Shiu applies because every variable exceeds $H=K^{10}$.  At the lowest dyadic interval, $z\geq X^{1/12}>K^{20}$ follows from $\kappa<1/240$.  Lemma~\ref{lem:incidence} and Markov's inequality bound the harmonic mass of discarded triples by
\[
 \frac{(\log z)^2(1+\log K)^3}{(\log X)^4}
 =o((\log z)^2h(\mathcal J)).
\]
Thus the retained first moment has full order.  For fixed $4uv$, coprimality permits at most $2^{\omega(uv)}$ allocations to $(u,v)$, and the pruning permits at most $(\log X)^4$ choices of $k$.  The retained progression multiplicity is therefore
$(\log X)^{D\log2+4}$.  Ordinary Bombieri--Vinogradov absorbs it.  Distinctness is unchanged, and Brun--Titchmarsh gives the matching upper bound.  Finally
$h(\mathcal K(X^\kappa))\asymp\log X$.
\end{proof}
Theorem~\ref{thm:pruned} does not prove the literal weighted error-sum hypothesis from which it originated, does not assemble composite moduli, and does not improve Theorem~\ref{thm:record} unconditionally.  Its soft spots are the power-range extension of the Shiu/box estimates and the uniform second-moment boundary sum; these require external review.

\section{The fixed-$c$-free critical-window assembly}\label{sec:cfree}
\begin{center}
\fbox{\parbox{0.91\linewidth}{\centering
\textsc{Claimed/provisional.}  Every result in this section which feeds
Theorem~\ref{thm:three-quarter} is subject to hostile external checking.
In particular, the load-bearing prime-slice input Theorem~\ref{thm:pruned}
is likewise provisional.  The argument is not presented as an established
record.}}
\end{center}

This section attacks the restricted critical-window assembly.  It does
\emph{not} prove the literal coefficient-sum clause of the source-notes
hypothesis $H_{\rm PF}'(k\ell;{\rm good})$.  Instead it removes the moving
$c$-family, uses one fixed unpruned atom family, and proves the weaker ledger
$\exp\{O(J\log X)\}$.  That ledger is of the same order as the lcm budget and
is still absorbed when $\log N\geq C(\log X)^4$.  Subject to the caveats in
Section~\ref{sec:verification-remarks}, this gives the
\textsc{claimed/provisional} unconditional $3/4$ exponent.

Put
\begin{equation}\tag{39.1}
 t=\log X,\qquad K=\lfloor X^\kappa\rfloor,
 \qquad H=K^{10},\qquad 0<\kappa<1/240,
\end{equation}
and partition $(X^{1/2},X]$ into disjoint dyadic blocks $(x,2x]$, truncating
the two endpoint blocks.  In a block put $z=x^{1/6}$.  Let $\mathcal A_X$
be the following fixed \emph{$c$-free} family.  An atom is
\begin{equation}\tag{39.2}
 A=(k,\ell,u,v),\quad k\leq K,\ k\equiv1\pmod4,
 \quad \ell\in(x,2x]\text{ prime},\quad H<u,v\leq z,
\end{equation}
with
\begin{equation}\tag{39.3}
 (u,v)=(uv,k)=1,\qquad
 \omega(uv)\leq D\log\log X,\qquad 4uv\mid k\ell+1.
\end{equation}
Here $D$ is the fixed low-$\omega$ cutoff constant used in the source Lemma
16.2 and in the proof of Theorem~\ref{thm:pruned}.  The atom denotes the
cylinder
\begin{equation}\tag{39.4}
 E_A=\{n:n\equiv-uv^{-1}\pmod{k\ell}\}.
\end{equation}
The dyadic cutoff is only a canonical subfamily of the multiplier-prime
family.  The lower proof of Theorem~\ref{thm:pruned}, specifically the
source-notes estimates (34.20)--(34.21), furnishes its retained main mass
inside these boxes.

\subsection{Removing the moving fibre and deduplicating the atoms}
\begin{lemma}[source Lemma 39.1: $c$-free reduction and distinct large
coordinates; proved, provisionally checked]\label{lem:cfree}
The following hold.
\begin{enumerate}[leftmargin=*]
\item If $n$ is exceptional, then it avoids every atom in $\mathcal A_X$.
If additionally $c\equiv n\pmod{24L_K}$, then every atom hit by $n$
automatically obeys $k\mid u+cv$.  Thus the coupling condition used to
define the source family $\mathscr G_{X,c}$ is implied by the atom itself;
it need not be built into the polynomial.
\item Distinct atoms which have a compatible intersection have distinct
primes $\ell$.  At a fixed $\ell$, all atom classes modulo $\ell$ are
distinct.  At a fixed modulus $k\ell$, no two ordered pairs $(u,v)$ produce
the same class.
\item Consequently, for a compatible ordered set $A_1,\ldots,A_j$,
\begin{equation}\tag{39.5}
 \Pr\Big(\bigcap_{i\leq j}E_{A_i}\Big)
 =\frac{\mathbf1_{\rm compatible}}
 {\lcm(k_1,\ldots,k_j)\prod_{i\leq j}\ell_i}.
\end{equation}
\end{enumerate}
\end{lemma}
\begin{proof}
Lemma~\ref{lem:multiplier} proves the first assertion.  Modulo $k$,
(39.4) says $u+nv\equiv0$; since $c\equiv n\pmod k$, it also says
$u+cv\equiv0$.

Suppose two atoms at the same $\ell$ are compatible.  Their projections
modulo $\ell$ agree, so
\[
 \ell\mid uv'-u'v.
\]
In the same block the absolute value is at most $z^2=x^{1/3}<\ell$.  The
same conclusion holds if two admissible cutoffs are compared across blocks,
since every $u,v\leq X^{1/6}$, while $\ell>X^{1/2}>X^{1/3}$.  Hence
$uv'=u'v$, and reducedness gives $(u,v)=(u',v')$.  The apparent swap can
agree only when $u=v$, which is impossible here because $(u,v)=1$ and
$u,v>H$.  Finally $uv\mid(k\ell+1)/4$ is exactly
\[
 k\ell\equiv-1\pmod{4uv}.
\]
Since $(\ell,4uv)=1$, this determines $k\pmod{4uv}$, and
$4uv>4H^2>K$ makes the admissible $k$ unique.  This proves both forms of
deduplication and the distinct-$\ell$ assertion.  The Chinese remainder
theorem now gives (39.5).
\end{proof}

The logical point is as important as the geometry.  The majorant below
covers the fixed predicate
$\mathbf1_{(n,P_y)=1}\mathbf1_{H_X(n)=0}$, where
$H_X=\sum_{A\in\mathcal A_X}\mathbf1_{E_A}$.  It does not majorize every
avoider of the smaller, $c$-dependently pruned source family.  It does
majorize every exceptional prime, which is all the exceptional-set
implication needs.

\subsection{Exact mass profiles}
For $g\mid k$ and $a\in(\mathbb Z/g\mathbb Z)^\times$, write
\begin{equation}\tag{39.6}
 W_{k,a}(g)=
 \sum_{\substack{A\in\mathcal A_X\text{ with multiplier }k\\
                  -uv^{-1}\equiv a\ (g)}}\frac1{k\ell},
 \qquad W_k=\sum_{A:k_A=k}\frac1{k\ell}.
\end{equation}

\begin{lemma}[source Lemma 39.2: multiplier and residue profiles; proved,
provisional]\label{lem:profiles}
Uniformly for $k\leq K$, $k\equiv1\pmod4$, $g\mid k$, and reduced
$a\pmod g$,
\begin{equation}\tag{39.7}
 W_{k,a}(g)\ll \frac{t^2}{\varphi(g)}\frac{\varphi(k)^2}{k^3},
 \qquad W_k\asymp t^2\frac{\varphi(k)^2}{k^3}.
\end{equation}
Consequently
\begin{equation}\tag{39.8}
 \mu_X:=\sum_{A\in\mathcal A_X}\frac1{k_A\ell_A}
 \asymp t^2\sum_{\substack{k\leq K\\k\equiv1(4)}}
             \frac{\varphi(k)^2}{k^3}
 \asymp t^2\log K\asymp t^3.
\end{equation}
For each fixed odd prime $p$, the share of the proxy Euler weight
$b(k)/k=\varphi(k)^2/k^3$ carried by multipliers divisible by $p$ is
exactly
\begin{equation}\tag{39.9}
 \alpha_p=\frac{p-1}{p^2+p-1}\sim\frac1p.
\end{equation}
Because (39.7) determines $W_k$ only up to constants, (39.9) is a statement
about the proxy weight, not an exact share of the atom mass; the moment proof
below never uses $\alpha_p$.
\end{lemma}
\begin{proof}
We give the review-inserted derivation rather than appeal to an average over
$c$.  Fix the block and $k$, and put $F(n)=n/\varphi(n)$ and
\[
 q_0=\lcm(g,\rad k)\leq k.
\]
Work in multiplicative boxes $u\asymp U$, $v\asymp V$.  The conditions
$(uv,k)=1$ and $-uv^{-1}\equiv a\pmod g$ restrict
$(u\bmod q_0,v\bmod q_0)$ to reduced pairs with a fixed ratio modulo $g$:
for each of the $\varphi(q_0)$ choices of $v$-class, the $u$-class is
determined modulo $g$ and free at the remaining part, so exactly
$\varphi(q_0)^2/\varphi(g)$ ordered class pairs are compatible.  Shiu's
theorem for the multiplicative function $F$, on the fixed-relative-length
interval $u\asymp U$ in one reduced class modulo $q_0$, is applicable since
$q_0\leq k\leq K$ and $U>H=K^{10}$, and gives
\[
 \sum_{\substack{u\asymp U\\u\equiv r\ (q_0)}}F(u)
 \ll \frac{U}{\varphi(q_0)}
 \exp\Big\{-\sum_{p\mid k}\frac1{p-1}\Big\}
 \asymp \frac{U}{\varphi(q_0)}\frac{\varphi(k)}k.
\]
The last comparison holds because the logarithm of the quotient is
$O(\sum_p p^{-2})$.  Multiplying the $u$- and $v$-class sums by the number
of compatible class pairs and dividing by $UV$ yields, per box pair,
\[
 \sum_{\rm box}\frac1{\varphi(u)\varphi(v)}
 \ll\frac1{\varphi(g)}\left(\frac{\varphi(k)}k\right)^2.
\]
Summing the $O((\log(z/H))^2)=O(t^2)$ box pairs gives
\begin{equation}\tag{39.11}
 \sum\frac1{\varphi(u)\varphi(v)}
 \ll\frac{(\log z)^2}{\varphi(g)}\frac{\varphi(k)^2}{k^2}.
\end{equation}
For fixed $(u,v)$, Brun--Titchmarsh in the progression
$\ell\equiv-k^{-1}\pmod{4uv}$, whose modulus is
$4uv\leq4x^{1/3}$, gives
\[
 \sum_{\substack{x<\ell\leq2x\\
                  \ell\equiv-k^{-1}\ (4uv)}}\frac1\ell
 \ll\frac1{\varphi(4uv)\log x},
 \qquad
 \frac1{\varphi(4uv)}\ll\frac1{\varphi(u)\varphi(v)},
\]
where the last inequality uses $(u,v)=1$.  Combining this with (39.11)
over the boxes proves
\begin{equation}\tag{39.10}
 \sum_{\substack{A\text{ in the block}\atop
       k_A=k,\ -uv^{-1}\equiv a\ (g)}}\frac1\ell
 \ll\frac{\log x}{\varphi(g)}\frac{\varphi(k)^2}{k^2},
\end{equation}
with an absolute constant, uniformly in $k\leq K$, $g\mid k$, and reduced
$a$: no $k^\varepsilon$, divisor, or logarithmic loss enters.  After the
$O(t)$ blocks and division by $k$, this is the upper half of (39.7).

For the lower half drop the ratio condition.  For one fixed $k$, a modulus
$4uv$ has at most $2^{\omega(uv)}=(\log X)^{O(1)}$ ordered allocations.
Thus ordinary Bombieri--Vinogradov, with a fixed sufficiently large
logarithmic saving, evaluates the prime progressions without the factor
$K$: $k$ is fixed on this line.  The low-$\omega$ Rankin argument from the
source proof retains a fixed proportion uniformly.  At a prime $p\mid k$
the exact coprime-pair local factor differs from $(1-1/p)^2$ by
$1+O(p^{-2})$, so the unrestricted harmonic pair mass is
$\asymp(\log(z/H))^2(\varphi(k)/k)^2$ uniformly.  This gives the matching
lower bound block by block and proves (39.7).

Finally $b(k)=(\varphi(k)/k)^2$ is multiplicative and
$\sum b(k)/k\asymp\log K$, also after selecting $k\equiv1\pmod4$, by the
principal/$\chi_4$ decomposition.  At an odd prime its local factor is
\begin{equation}\tag{39.12}
 1+\sum_{e\geq1}\frac{(1-1/p)^2}{p^e}=1+\frac{p-1}{p^2},
\end{equation}
so the part with $p\mid k$, divided by the full factor, is exactly (39.9).
\end{proof}

Given a reduced fibre $c$, an atom with multiplier $k$ is active for one of
the $\varphi(k)$ unit residues and then has prime-coordinate probability
$1/\ell$.  Averaging those conditional masses produces the factor
$1/\varphi(k)$, whereas the unconditioned cylinder has probability
$1/(k\ell)$.  Formula (39.7), not the conditional mass profile used in the
earlier moving-$c$ family, is the correct profile for factorial moments in
the $c$-free CRT space.

\subsection{Collapse versus consistency}
The rough collapse bound which discards residue consistency would charge a
shared prime $p$ by $p\alpha_p^2\asymp1/p$, producing a false
$\log\log K$ loss.  The required cancellation is exact at the level of a
new compatible atom.

For a cutoff $y$, put
\begin{equation}\tag{39.13}
 q_y(k)=\prod_{\substack{p\mid k\\p\leq y}}\frac p{p-1},
 \qquad
 b_y(g)=\prod_{\substack{p^e\Vert g\\p\leq y}}\varphi(p^e)
         \prod_{\substack{p^e\Vert g\\p>y}}p^e.
\end{equation}
Condition the CRT space by $(n,P_y)=1$, where $P_y=\prod_{p\leq y}p$.
An atom then has probability $q_y(k)/(k\ell)$.  If a compatible previously
selected set has multiplier lcm $R$, adding an atom of multiplier $k$, with
$g=(k,R)$, multiplies its intersection probability by
\begin{equation}\tag{39.14}
 \frac{q_y(k)}{k\ell}\,b_y(g).
\end{equation}
Compatibility fixes one unit residue modulo $g$.  Lemma~\ref{lem:profiles}
says that the atom mass in that residue is at most $1/\varphi(g)$ of its
multiplier mass.  Thus the collapse left after summing the compatible new
atoms is
\begin{equation}\tag{39.15}
 \frac{b_y(g)}{\varphi(g)}
 =\prod_{\substack{p\mid g\\p>y}}\frac p{p-1};
\end{equation}
small conditioned primes cancel exactly, and an unconditioned shared prime
costs only $p/(p-1)$, not $p$.

\begin{lemma}[source Lemma 39.3: uniform multiplier Euler bound; proved]
\label{lem:euler-bound}
Uniformly in $y\geq2$, every integer $R$, and $K\geq3$,
\begin{equation}\tag{39.16}
 \sum_{\substack{k\leq K\\k\equiv1(4)}}
 \frac{\varphi(k)^2}{k^3}q_y(k)
 \prod_{\substack{p\mid(k,R)\\p>y}}\frac p{p-1}
 \ll\log K.
\end{equation}
\end{lemma}
\begin{proof}
It is enough to sum over all odd $k$.  The unmodified local nonconstant mass
is $(p-1)/p^2$.  Multiplication either by $q_y$ for $p\leq y$, or by
$p/(p-1)$ when $p>y$ and $p\mid R$, changes this to $1/p$.  The quotient
\begin{equation}\tag{39.17}
 \frac{1+1/p}{1+(p-1)/p^2}=1+O(p^{-2})
\end{equation}
has a convergent product, uniformly in the chosen set of modified primes.
The remaining harmonic Euler product is $O(\log K)$.
\end{proof}

\begin{theorem}[source Theorem 39.4: all required upper factorial moments;
proved, provisional]\label{thm:factorial-cfree}
In the CRT space conditioned by $(n,P_y)=1$, for every integer $m\geq1$,
\begin{equation}\tag{39.18}
 \mathbb E(H_X)_m\leq(Ct^3)^m,
\end{equation}
with $C$ depending at most on $\kappa$, uniformly in
$2\leq y<X^{1/2}$.  The upper restriction is necessary: for $y\geq\ell$
the coprimality conditioning would renormalize the $\ell$-coordinate itself
by $\ell/(\ell-1)$, which (39.14) does not include.  The application below
has $y=Bt^3<X^{1/2}$ for all sufficiently large $X$.  In particular the
bound holds through every order $m=D_Bt^3$.
\end{theorem}
\begin{proof}
Expand the ordered factorial moment into distinct compatible atoms.  By
Lemma~\ref{lem:cfree}, every admissible next atom has a new $\ell$.  For a
fixed preceding compatible set, (39.14), the residue-profile upper bound
(39.7), and (39.15) bound the sum over all possible next atoms by
\begin{equation}\tag{39.19}
 Ct^2\sum_{\substack{k\leq K\\k\equiv1(4)}}
 \frac{\varphi(k)^2}{k^3}q_y(k)
 \prod_{\substack{p\mid(k,R)\\p>y}}\frac p{p-1}
 \leq Ct^2\log K\leq Ct^3.
\end{equation}
Removing the forbidden old $\ell$'s only decreases this sum.  Iteration
proves (39.18).
\end{proof}

At pair level, for fixed $k,k'$ and $g=(k,k')$,
\begin{equation}\tag{39.20}
 \sum_{A:k_A=k}\sum_{\substack{B:k_B=k'\\A\sim B}}
 \Pr(E_A\cap E_B)
 \ll\frac g{\varphi(g)}W_kW_{k'}.
\end{equation}
Averaging (39.20) uses
$\alpha_p^2(p/(p-1)-1)=O(p^{-3})$, not
$p\alpha_p^2\asymp1/p$.  Hence the pair excess is $O(\mu_X^2)$ with an
absolute constant.  It is \emph{not} $o(\mu_X)$: a fixed small prime already
contributes a constant multiple of $\mu_X^2$ to the usual Janson dependency
sum.  Therefore the proposed direct condition $\Delta=o(\mu)$ and the naive
Janson proof fail quantitatively.  The factorial-moment theorem survives
because an absolute base $C^m$, rather than Poisson relative error, is enough
after the separate void estimate below.

\subsection{The CRT void without partitioning by all of $M_0$}
A completely unconditioned $c$-free void still has a quarantine floor:
taking $n\equiv0\pmod p$ for many small $p$ disables every atom whose
multiplier contains one of them.  Thus $Q_r(H_X)$ alone cannot have cubic
mean.  Take
\begin{equation}\tag{39.21}
 y=Bt^3
\end{equation}
with a sufficiently large fixed $B$, and multiply by the exact small-prime
coprimality selector
\begin{equation}\tag{39.22}
 S_y(n)=\mathbf1_{(n,P_y)=1}
 =\sum_{d\mid P_y}\mu(d)\mathbf1_{d\mid n}.
\end{equation}
Every exceptional prime greater than $y$ has $S_y=1$.  Notice that (39.22)
has degree $\pi(y)=O(t^3/\log t)$, not the impossible degree $\pi(K)$.

\begin{lemma}[source Lemma 39.5: conditioned $c$-free void; proved,
provisional]\label{lem:conditioned-void}
In the exact CRT space,
\begin{equation}\tag{39.23}
 \Pr(H_X=0\mid S_y=1)\leq e^{-ct^3}.
\end{equation}
\end{lemma}
\begin{proof}
Reveal all multiplier coordinates.  Let $c$ denote the resulting class
modulo $24L_K$, where
$L_K=\lcm\{k\leq K:k\equiv1\pmod4\}$, and put
\begin{equation}\tag{39.24}
 \mathcal J_c=\{k\leq K:k\equiv1\pmod4,\ (k,c)=1\},
 \qquad
 Z(c)=\sum_{\substack{y<p\leq K,\ p\mid L_K\\p\mid c}}\frac1p.
\end{equation}
The restriction $p\mid L_K$ makes $Z$ well defined: a prime
$p\equiv3\pmod4$ with $3p>K$ divides no admissible multiplier, so it has
no CRT coordinate here and can be omitted; it removes nothing from
$\mathcal J_c$.  The coordinates for $p>y$, $p\mid L_K$, remain independent
and $\Pr(p\mid c)=1/p$.  With exponential parameter $y$,
\begin{equation}\tag{39.25}
 \mathbb E e^{yZ}
 \leq\exp\left\{Cy\sum_{p>y}p^{-2}\right\}=e^{O(1)},
 \qquad \Pr(Z>\eta)\leq e^{-\eta y+O(1)}.
\end{equation}

Write $h(\mathcal J)=\sum_{k\in\mathcal J}\varphi(k)/k^2$.  The elementary
bounds
\begin{equation}\tag{39.26}
 h(\mathcal K(K))\asymp\log K,
 \qquad
 \sum_{\substack{k\leq K\\p\mid k}}\frac{\varphi(k)}{k^2}
 \ll\frac{\log K}{p}
\end{equation}
show by the union bound that, for a sufficiently small fixed $\eta$,
$Z(c)\leq\eta$ implies
$h(\mathcal J_c)\geq c_1\log K$.  Also $S_y=1$ makes $c$ coprime to $24$,
and by definition it is coprime to $L_{\mathcal J_c}$.

Apply Theorem~\ref{thm:pruned}, the source-notes Theorem 34.8, to this
arbitrary subfamily $\mathcal J_c$.  Its lower proof lies in the canonical
boxes (39.2), and its good triples are contained in the $c$-free unpruned
family.  Conditional on the multiplier coordinates, the distinct
$\ell$-coordinates are independent by Lemma~\ref{lem:cfree}.  Therefore,
when $Z(c)\leq\eta$,
\begin{equation}\tag{39.27}
 \Pr(H_X=0\mid c)
 \leq\prod_\ell\left(1-\frac{f_c^{\rm good}(\ell)}\ell\right)
 \leq\exp\{-c_2t^2h(\mathcal J_c)\}
 \leq e^{-c_3t^3}.
\end{equation}
The bad multiplier coordinates cost at most
$e^{-\eta Bt^3+O(1)}$ by (39.25).  Increasing $B$ if necessary and averaging
(39.27) proves (39.23).
\end{proof}

This is the load-bearing use of Theorem~\ref{thm:pruned}.  It replaces
Janson by supplying cubic active prime-coordinate mass uniformly for every
reduced surviving multiplier fibre.  The only fibres not reduced are those
which set too much large-prime multiplier mass to zero, and (39.25) makes
those fibres exponentially rare after the polynomial-sized coprimality
cutoff.  The dependency on that provisional theorem is therefore explicit
and unavoidable in the present argument.

\subsection{A low-degree majorant with the honest ledger}
For even $r$, put
\[
 Q_r(h)=\sum_{j=0}^r(-1)^j\binom hj.
\]
The exact identity
\begin{equation}\tag{39.28}
 Q_r(0)=1,\qquad Q_r(h)=\binom{h-1}r\geq0\quad(h\geq1)
\end{equation}
shows that
\begin{equation}\tag{39.29}
 \nu_X(n)=S_y(n)Q_r(H_X(n))
\end{equation}
is nonnegative and is at least one on every exceptional prime
$n>\max(K,y)$.

\begin{theorem}[source Theorem 39.6: modified $c$-free critical-window
assembly; proved, \textsc{claimed/provisional}]\label{thm:cfree-assembly}
There are constants $c,C_0,D_B>0$, depending at most on $\kappa$, such that
if $r$ is the least even integer at least $D_Bt^3$, then (39.29) has exact
CRT mean
\begin{equation}\tag{39.30}
 \mathbb E_{\rm CRT}\nu_X\leq e^{-ct^3}.
\end{equation}
Here and below $D_B$ is the Bonferroni-degree constant, distinct from the
source Lemma 16.2's $\omega$-cutoff constant $D$ in (39.3).  Its congruence
expansion has the explicit bounds
\begin{equation}\tag{39.31}
 \begin{split}
 \text{degree}&\leq r+\pi(y)=O(t^3),\\
 \log d_{\rm term}&\leq O(y)+r(t+\log K)=O(t^4),\\
 \log\sum_{\rm terms}|c_{\rm term}|&=O(rt)=O(t^4).
 \end{split}
\end{equation}
Every nonzero term is one plain congruence class.  Consequently, whenever
$\log N\geq C_0t^4$,
\begin{equation}\tag{39.32}
 \sum_{n\leq N}\nu_X(n)\ll Ne^{-ct^3}.
\end{equation}
\end{theorem}
\begin{proof}
From (39.28), for $h\geq1$,
$\binom{h-1}r\leq\binom h{r+1}$.  Lemma~\ref{lem:conditioned-void} and
Theorem~\ref{thm:factorial-cfree} therefore give
\begin{equation}\tag{39.33}
 \begin{split}
 \mathbb E(Q_r(H_X)\mid S_y=1)
 &\leq\Pr(H_X=0\mid S_y=1)
       +\frac{\mathbb E((H_X)_{r+1}\mid S_y=1)}{(r+1)!}\\
 &\leq e^{-c_1t^3}+\frac{(Ct^3)^{r+1}}{(r+1)!}
 \leq e^{-ct^3},
 \end{split}
\end{equation}
when $D_B$ is sufficiently large.  Multiplication by
$\Pr(S_y=1)\leq1$ proves (39.30).

Expand (39.22) and the elementary symmetric polynomials in $Q_r$.
Lemma~\ref{lem:cfree} makes every compatible atom set use distinct
$\ell$'s; CRT then makes its further intersection with $d\mid P_y$ either
empty or one class.  If $A_X=|\mathcal A_X|$, the crude choice count from
(39.2) gives $\log A_X=O(t)$, and hence
\begin{equation}\tag{39.34}
 2^{\pi(y)}\sum_{j\leq r}\binom{A_X}j=\exp\{O(rt)\}.
\end{equation}
The modulus is at most $P_y(KX)^r$, proving (39.31).  On $[1,N]$, each
class count is its CRT mean times $N$, plus $O(1)$.  Choosing $C_0$ larger
than the constants in (39.31) makes the aggregate rounding error
$\exp\{O(rt)\}$ at most, say, $N^{1/2}$, which is negligible beside
$Ne^{-ct^3}$.  This proves (39.32).
\end{proof}

The literal source hypothesis $H_{\rm PF}'(k\ell;{\rm good})$ demanded
coefficient sum $e^{O(J)}$, with $J\asymp t^3$, and a majorant for its larger
aggregate avoider predicate.  Theorem~\ref{thm:cfree-assembly} proves
\emph{neither} clause: its direct even-Bonferroni expansion has ledger
$e^{O(Jt)}$, and it only covers the $c$-free avoider, hence all exceptional
primes.  These are deliberate modifications, not suppressed losses.  The
ledger still closes because the critical window itself has logarithmic size
$\asymp Jt$.  Thus the literal $e^{O(J)}$ clause remains \textsc{open} as
stated, but is no longer needed for the exceptional set if
Theorem~\ref{thm:cfree-assembly} survives review.

\subsection{Exceptional denominators}
\begin{theorem}[source Theorem 39.7: unconditional $3/4$ exceptional-set
bound; \textsc{claimed/provisional}]\label{thm:three-quarter}
There is a constant $c>0$ such that
\begin{equation}\tag{39.35}
 E_{\rm all}(N)\ll N\exp\{-c(\log N)^{3/4}\}.
\end{equation}
The same bound holds for exceptional primes.
\end{theorem}
\begin{proof}[Internal argument; provisional status retained]
Put $L=\log N$, take $t=\alpha L^{1/4}$, and set $X=e^t$.  Choose the fixed
$\alpha>0$ small enough that $L\geq C_0t^4$.
Lemma~\ref{lem:cfree} and (39.29) majorize every exceptional prime larger
than $\max(K,y)$; that omitted range is negligible.  Equations (39.32) and
$t^3=\alpha^3L^{3/4}$ give the prime bound.

For completeness, we spell out the semigroup transfer used for all
denominators.  Every prime factor of an exceptional integer $n>1$ is itself
exceptional: if $p\mid n$ and
$4/p=1/x_1+1/x_2+1/x_3$, scaling the three denominators by $n/p$ represents
$4/n$.  Thus the exceptional integers are contained in the multiplicative
semigroup generated by exceptional primes, including the empty product.
The integer $1$ contributes only $O(1)$.

Write $g(u)=u^{3/4}$ and $x=e^L$.  Fix $x_0$ large enough that the prime
bound gives, for every $x\geq x_0$,
\[
 E_{\rm prime}(x)\ll x\exp\{-c_0g(\log x)\}.
\]
Fix $0<\eta<c_0$ and, for large $L$, put
$\delta=\eta g(L)/L$, so $0<\delta<1/4$.  Rankin's inequality over the
semigroup gives
\[
 E_{\rm all}(x)\leq x^{1-\delta}
 \prod_{\substack{p\leq x\\p\ \text{exceptional}}}
       (1-p^{-1+\delta})^{-1}.
\]
The finitely many exceptional primes $p\leq x_0$ contribute $O(1)$,
uniformly in $\delta$.  Partial summation and the endpoint estimate
$\delta L-c_0g(L)=-(c_0-\eta)g(L)<0$ give
\[
 \sum_{\substack{p\leq x\\p\ \text{exceptional}}}p^{-1+\delta}
 \ll1+\int_{\log x_0}^L\exp\{\delta u-c_0g(u)\}\,du.
\]
Because $g(u)/u=u^{-1/4}$ is decreasing, for $u\leq L$ one has
$\delta u\leq\eta g(u)$.  The last integral is therefore bounded by
$\int_{\log x_0}^\infty e^{-(c_0-\eta)g(u)}\,du<\infty$.  The prime-power
terms in the logarithm of the Euler product also satisfy
\[
 \sum_{p\leq x}\sum_{j\geq2}\frac{p^{-j(1-\delta)}}j
 \leq\sum_p\sum_{j\geq2}\frac{p^{-3j/4}}j=O(1).
\]
Hence the Euler product is $O(1)$ and
$E_{\rm all}(x)\ll x^{1-\delta}=x\exp\{-\eta g(\log x)\}$, proving
(39.35).
\end{proof}

\begin{center}
\fbox{\parbox{0.91\linewidth}{\centering
Theorem~\ref{thm:three-quarter} is \textsc{claimed/provisional}.  It uses
the likewise provisional Theorem~\ref{thm:pruned}.  It does not use the
source hypotheses $H_{k\rm BV}$, $H_{\rm PF}'$, or Elliott--Halberstam, but
``unconditional'' does not remove the need for external verification.}}
\end{center}

\subsection{Remarks on verification status}\label{sec:verification-remarks}
The three weakest steps identified by the maximum-severity self-review are
reproduced faithfully here.
\begin{enumerate}[leftmargin=*]
\item The new residue-resolved upper bound (39.11) must be rederived by a
referee with every local factor present.  It needs simultaneous uniformity
in $k\leq X^\kappa$, every divisor $g\mid k$, and every unit residue; the
proof uses Brun--Titchmarsh for the upper bound, so it does not hide a
polynomially small Bombieri--Vinogradov main term.  This is the most
load-bearing new analytic lemma.
\item Lemma~\ref{lem:conditioned-void} uses Theorem~\ref{thm:pruned} for the
data-dependent subfamily $\mathcal J_c$ and uses the canonical lower boxes,
not merely the larger family in the statement of the source formula
(34.18).  The quantifiers of Theorem~\ref{thm:pruned} do allow every
subfamily containing $1$, and its proof furnishes the lower mass in those
boxes, but this inheritance is new and needs an independent line-by-line
check together with Theorem~\ref{thm:pruned} itself.
\item The ordered-moment induction (39.14)--(39.19) is where lcm collapse
and residue consistency cancel at every prime power.  A missed factor there
would be exponentiated through $r\asymp t^3$.  The displayed conditional
factor is $\varphi(p^e)$ for $p\leq y$ and $p^e$ for $p>y$; after the
$1/\varphi(g)$ residue share, only $p/(p-1)$ remains at an unconditioned
shared prime.  This exact replay, including unequal prime-power exponents,
is mandatory in external review.
\end{enumerate}

Internally, the quantifier chain is complete: Theorem~\ref{thm:pruned} is
uniform in $\mathcal J_c$; (39.25) removes the nonreduced fibres; (39.18)
controls the even Bonferroni tail; and (39.31)--(39.34) pay the full
$e^{O(t^4)}$ coefficient ledger before selecting
$t=\alpha(\log N)^{1/4}$.  The direct Janson proposal is rejected because
its dependency sum is $\Theta(\mu^2)$, while the original $e^{O(J)}$
formulation remains open.  Theorem~\ref{thm:three-quarter} should not be
cited as established outside this campaign until an independent expert has
checked the three steps above and the priority claim against the primary
literature.

A fresh search found no post-1970 improvement of Vaughan's
$N/\exp(c(\log N)^{2/3})$ exceptional-set bound.  Pomerance--Weingartner's
2025 preprint \cite[\S4]{PomeranceWeingartner} explicitly cites Vaughan
1970 as the state of the art and reproduces the same $2/3$ exponent when
its general theorem is specialized to $m=4$.  Access to Vaughan's paper was
blocked, but DOI \href{https://doi.org/10.1112/S0025579300002886}
{10.1112/S0025579300002886} was confirmed.  If
Theorem~\ref{thm:three-quarter} survives external review, it would therefore
be the first exponent improvement found since 1970.  This is a search
result, not a literature guarantee; the \textsc{claimed/provisional} label
stands.

\subsection{Residue-dispersion frontier}\label{sec:dispersion-frontier}
Put
\[
 L=\log X,\qquad z=\frac{L^3\log\log L}{\log L},\qquad
 \Lambda=\frac{L^3}{\log L}.
\]
For a composite $z$-rough $M\leq X$, $M\equiv3\pmod4$, choose the least
positive divisor $D\mid((M+1)/4)^2$ in each fibre of
$D\mapsto-4D\pmod M$, and delete a representative when its class is implied
by a prime-modulus atom.  Call the remaining set $\mathcal D^\circ(M)$.
With the exact conditional weight
\[
 \kappa(M)=\prod_{\substack{q\mid M,\ z<q\leq Y\\q\equiv3(4)}}
                \frac q{q-F(q)},
 \qquad
 n_{M,p}(a)=\#\{D\in\mathcal D^\circ(M):-4D\equiv a\pmod p\},
\]
the pair-level dispersion quantity is exactly
\begin{equation}\tag{40.5}
 \mathcal V_X=\sum_{z<p\leq X}p\sum_{a\bmod p}
 \left(\sum_{\substack{M\leq X,
      \ M\ \text{composite},\ P^-(M)>z\\
      M\equiv3(4),\ p\mid M}}
      \frac{\kappa(M)}M n_{M,p}(a)\right)^2.
\end{equation}
This is the residue-dispersion reformulation; canonical deduplication,
prime powers, and the prime-implied deletion are part of the formula.

Write $D=R^2/s$, with $s\mid\rad(R)$ and $M=4Rk-1$.  For
$p\nmid4R$, let $c_p(R)\in\{1,\ldots,p-1\}$ be the inverse of $4R$ modulo
$p$.  The condition $p\mid M$ gives $k=c_p(R)+jp$.  Splitting into the
endpoint $j=0$ and regular tail $j\geq1$ gives vectors
$u^{\rm end}_{p,a}$ and $u^{\rm reg}_{p,a}$.  The regular-tail theorem is
proved, with the full canonical and retention restrictions allowed:
\begin{equation}\tag{40.14}
 \mathcal V_X^{\rm reg}
 :=\sum_{z<p\leq X}p\sum_a(u^{\rm reg}_{p,a})^2
 \ll L^2\log L+\frac{L^7}{z\log z}=o(\Lambda^2).
\end{equation}
Indeed the contribution of fixed $D$ is at most $CL/(pR_0(D))$; after
expansion, $p\mid D-E$, while
$\sum_D L/R_0(D)\ll L^3$ and
$\sum_D(L/R_0(D))^2\ll L^2$.  For $D\ne E$ there are at most
$2L/\log z$ prime divisors above $z$, and for $D=E$ Mertens' estimate costs
$O(\log L)$, giving (40.14).

The exact remaining pair target is therefore
\begin{equation}\tag{40.18}
 \boxed{\quad \mathcal V_X^{\rm end}
 =\sum_{z<p\leq X}p\sum_a(u^{\rm end}_{p,a})^2
 =O(\Lambda^2).\quad}
\end{equation}
Equivalently, in divisor variables, it is the explicit short-cofactor
estimate
\begin{equation}\tag{40.19}
\begin{aligned}
 \sum_{z<p\leq X}p\sum_{a\bmod p}
 \Bigg(&\sum_{\substack{R,s,c,q:\ s\mid\rad(R),\ 1\leq c<p\\
       pq=4Rc-1\leq X,\ P^-(pq)>z,\ pq\equiv3(4)\\
       R^2/s\equiv-4^{-1}a\ (p),\ q>1\ ({\rm so}\ pq\ \text{composite})\\
       R^2/s\ \text{is\ the\ retained\ representative\ at}\ pq}}
       \frac{\kappa(pq)}{pq}\Bigg)^2
 \ll\Lambda^2.
\end{aligned}
\end{equation}
Every term has $q>z$ and $q<4R$.  Estimate (40.19) is \textsc{open}.
Standard additive large-sieve and ordinary divisor
Barban--Davenport--Halberstam statements do not contain this moving
inverse/short-cofactor family.  Even a proof of (40.19) would pass only the
pair test: the prime-power and multi-coordinate codegrees through
$m\asymp\Lambda$, the needed factorial-moment hierarchy, and the critical-
window avoider assertion all remain \textsc{open}.  No theorem about
$\PFp$ or the literal $H_{\rm PF}'$ follows from the proved regular tail.

\section*{Acknowledgment of status}
This is a consolidation draft, not a submission.  Computational observations from the research log are omitted unless needed to state a status fact.  The mathematical headline is the \textsc{claimed/provisional} $3/4$ exceptional-set bound, subject to external verification of its full chain and especially Theorem~\ref{thm:pruned}.  The natural-density refutation of the literal $\PF$ quantifiers is proved independently; the literal critical-window ledger clause and the short-cofactor dispersion endpoint remain open.

\begin{thebibliography}{99}

\bibitem{BloomElsholtz}
T.~F. Bloom and C. Elsholtz,
\emph{Egyptian fractions},
Nieuw Arch. Wiskd. (5) \textbf{23} (2022), no.~4, 237--245.

\bibitem{ElsholtzPlanitzer}
C. Elsholtz and S. Planitzer,
\emph{The number of solutions of the Erd\H{o}s--Straus equation and sums of $k$ unit fractions},
Proc. Roy. Soc. Edinburgh Sect. A \textbf{150} (2020), no.~3, 1401--1427;
arXiv:1805.02945.

\bibitem{ElsholtzTao}
C. Elsholtz and T. Tao,
\emph{Counting the number of solutions to the Erd\H{o}s--Straus equation on unit fractions},
J. Aust. Math. Soc. \textbf{94} (2013), 50--105;
arXiv:1107.1010.

\bibitem{Henriot}
K. Henriot,
\emph{Nair--Tenenbaum bounds uniform with respect to the discriminant},
Math. Proc. Cambridge Philos. Soc. \textbf{152} (2012), no.~3, 405--424;
arXiv:1102.1643.

\bibitem{NairTenenbaum}
M. Nair and G. Tenenbaum,
\emph{Short sums of certain arithmetic functions},
Acta Math. \textbf{180} (1998), 119--144.

% TODO(submission): update this draft citation if final publication metadata is used.
\bibitem{PomeranceWeingartner}
C. Pomerance and A. Weingartner,
\emph{Exceptions to the Erd\H{o}s--Straus--Schinzel conjecture},
author-hosted draft dated 20 November 2025, arXiv:2511.16817.

% TODO(submission): no journal publication was verified for this computational preprint.
\bibitem{Salez}
S.~E. Salez,
\emph{The Erd\H{o}s--Straus conjecture: new modular equations and checking up to $N=10^{17}$},
arXiv:1406.6307 (2014).

\bibitem{ScottSokal}
A.~D. Scott and A.~D. Sokal,
\emph{The repulsive lattice gas, the independent-set polynomial, and the Lov\'asz local lemma},
J. Stat. Phys. \textbf{118} (2005), 1151--1261 [TODO: journal volume/pages unverified from the archived PDF; arXiv:math/0309352 verified];
arXiv:cond-mat/0309352.

\bibitem{Shiu}
P. Shiu,
\emph{A Brun--Titchmarsh theorem for multiplicative functions},
J. Reine Angew. Math. \textbf{313} (1980), 161--170.

\bibitem{Vaughan}
R.~C. Vaughan,
\emph{On a problem of Erd\H{o}s, Straus and Schinzel},
Mathematika \textbf{17} (1970), 193--198,
DOI: \href{https://doi.org/10.1112/S0025579300002886}{10.1112/S0025579300002886}.

\bibitem{Yamamoto}
K. Yamamoto,
\emph{On the Diophantine equation $4/n=1/x+1/y+1/z$},
Mem. Fac. Sci. Kyushu Univ. Ser. A \textbf{19} (1965), no.~1, 37--47.

\end{thebibliography}
\end{document}
